%%% Hlavní soubor. Zde se definují základní parametry a odkazuje se na ostatní části. %%%

% Meta-data o práci (je nutno upravit)
\input metadata.tex

% Vygenerujeme metadata ve formátu XMP pro použití balíčkem pdfx
\input xmp.tex

%% Verze pro jednostranný tisk:
% Okraje: levý 40mm, pravý 25mm, horní a dolní 25mm
% (ale pozor, LaTeX si sám přidává 1in)
\documentclass[12pt,a4paper]{report}
\setlength\textwidth{145mm}
\setlength\textheight{247mm}
\setlength\oddsidemargin{15mm}
\setlength\evensidemargin{15mm}
\setlength\topmargin{0mm}
\setlength\headsep{0mm}
\setlength\headheight{0mm}
% \openright zařídí, aby následující text začínal na pravé straně knihy
\let\openright=\clearpage

%% Pokud tiskneme oboustranně:
% \documentclass[12pt,a4paper,twoside,openright]{report}
% \setlength\textwidth{145mm}
% \setlength\textheight{247mm}
% \setlength\oddsidemargin{14.2mm}
% \setlength\evensidemargin{0mm}
% \setlength\topmargin{0mm}
% \setlength\headsep{0mm}
% \setlength\headheight{0mm}
% \let\openright=\cleardoublepage

%% Pokud práci odevzdáváme pouze elektronicky, vypadají lépe symetrické okraje
% \documentclass[12pt,a4paper]{report}
% \setlength\textwidth{145mm}
% \setlength\textheight{247mm}
% \setlength\oddsidemargin{10mm}
% \setlength\evensidemargin{10mm}
% \setlength\topmargin{0mm}
% \setlength\headsep{0mm}
% \setlength\headheight{0mm}
% \let\openright=\clearpage

%% Vytváříme PDF/A-2u
\usepackage[a-2u]{pdfx}

%% Přepneme na českou sazbu a fonty Latin Modern
\usepackage[czech]{babel}
\usepackage{lmodern}

% Pokud nepouživáme LuaTeX, je potřeba ještě nastavit kódování znaků
\usepackage{iftex}
\ifpdftex
\usepackage[utf8]{inputenc}
\usepackage[T1]{fontenc}
\usepackage{textcomp}
\fi

%%% Další užitečné balíčky (jsou součástí běžných distribucí LaTeXu)
\usepackage{amsmath}        % rozšíření pro sazbu matematiky
\usepackage{amsfonts}       % matematické fonty
\usepackage{amsthm}         % sazba vět, definic apod.
\usepackage{bm}             % tučné symboly (příkaz \bm)
\usepackage{booktabs}       % lepší vodorovné linky v tabulkách
\usepackage{caption}        % umožní definovat vlastní popisky plovoucích objektů
\usepackage{csquotes}       % uvozovky závislé na jazyku
\usepackage{dcolumn}        % vylepšené zarovnání sloupců tabulek
\usepackage{floatrow}       % umožní definovat vlastní typy plovoucích objektů
\usepackage{graphicx}       % vkládání obrázků
\usepackage{icomma}         % inteligetní čárka v matematickém módu
\usepackage{indentfirst}    % zavede odsazení 1. odstavce kapitoly
\usepackage[nopatch=item]{microtype}  % mikrotypografická rozšíření
\usepackage{paralist}       % lepší enumerate a itemize
\usepackage[nottoc]{tocbibind} % zajistí přidání seznamu literatury,
                            % obrázků a tabulek do obsahu
\usepackage{xcolor}         % barevná sazba

% Balíček hyperref, kterým jdou vyrábět klikací odkazy v PDF,
% ale hlavně ho používáme k uložení metadat do PDF (včetně obsahu).
% Většinu nastavítek přednastaví balíček pdfx.
\hypersetup{unicode}
\hypersetup{breaklinks=true}

% Balíčky pro sazbu informatických prací
\usepackage{algpseudocode}  % součást balíčku algorithmicx
\usepackage[Algoritmus]{algorithm}
\usepackage{fancyvrb}       % vylepšené prostředí verbatim
\usepackage{listings}       % zvýrazňování syntaxe zdrojových textů

% Cleveref může zjednodušit odkazování, ale jeho užitečnost pro češtinu
% je minimalní, protože nezvládá skloňování.
% \usepackage{cleveref}

% Formátování bibliografie (odkazů na literaturu)
% Detailní nastavení můžete upravit v souboru macros.tex.
%
% POZOR: Zvyklosti různých oborů a kateder se liší. Konzultujte se svým
% vedoucím, jaký formát citací je pro vaši práci vhodný!
%
% Základní formát podle normy ISO 690 s číslovanými odkazy
\usepackage[natbib,style=iso-numeric,sorting=none]{biblatex}
% ISO 690 s alfanumerickými odkazy (zkratky jmen autorů)
%\usepackage[natbib,style=iso-alphabetic]{biblatex}
% ISO 690 s citacemi tvaru Autor (rok)
%\usepackage[natbib,style=iso-authoryear]{biblatex}
%
% V některých oborech je běžnější obyčejný formát s číslovanými odkazy
% (sorting=none říká, že se bibliografie má řadit podle pořadí citací):
%\usepackage[natbib,style=numeric,sorting=none]{biblatex}
% Číslované odkazy, navíc se [1,2,3,4,5] komprimuje na [1-5]
%\usepackage[natbib,style=numeric-comp,sorting=none]{biblatex}
% Obyčejný formát s alfanumerickými odkazy:
%\usepackage[natbib,style=alphabetic]{biblatex}

% Z tohoto souboru se načítají položky bibliografie
\addbibresource{literatura.bib}

% Definice různých užitečných maker (viz popis uvnitř souboru)
\input macros.tex

%%% Titulní strana a různé povinné informační strany
\begin{document}
%%% Titulní strana práce a další povinné informační strany

%%% Nápisy na přední straně desek
%%% Pokud je práce ve slovenštině, desky mají být česky.

% Desky obvykle nesázíme, ale pokud je chcete přidat, změnte \iffalse na \iftrue
\iffalse

\pagestyle{empty}
\hypersetup{pageanchor=false}
\begin{center}

\large
Univerzita Karlova

\medskip

Matematicko-fyzikální fakulta

\vfill

{\huge\bf\ThesisTypeTitle}

\vfill

{\huge\bf\ThesisTitle\par}

\vfill
\vfill

\hbox to \hsize{\YearSubmitted\hfil \ThesisAuthor}

\end{center}

\newpage\openright
\setcounter{page}{1}

\fi

%%% Titulní strana práce
%%% Pokud je práce ve slovenštině, tato strana zůstává česky.

\pagestyle{empty}
\hypersetup{pageanchor=false}

\begin{center}

\centerline{\mbox{\includegraphics[width=166mm]{img/logo-cs.pdf}}}

\vspace{-8mm}
\vfill

{\bf\Large\ThesisTypeTitle}

\vfill

{\LARGE\ThesisAuthor}

\vspace{15mm}

{\LARGE\bfseries\ThesisTitle\par}

\vfill

\Department

\vfill

{
\centerline{\vbox{\halign{\hbox to 0.45\hsize{\hfil #}&\hskip 0.5em\parbox[t]{0.45\hsize}{\raggedright #}\cr
Vedoucí \ThesisTypeGenitive{} práce:&\Supervisor \cr
\ifx\ThesisType\TypeRig\else
\noalign{\vspace{2mm}}
Studijní program:&\StudyProgramme \cr
\fi
}}}}

\vfill

Praha \YearSubmitted

\end{center}

\newpage

%%% Strana s čestným prohlášením k práci
%%% Pokud je práce ve slovenštině, tato strana zůstává česky.

\openright
\hypersetup{pageanchor=true}
\vglue 0pt plus 1fill

\noindent
Prohlašuji, že jsem tuto \ThesisTypeAccusative{} práci vypracoval(a) samostatně a výhradně
s~použitím citovaných pramenů, literatury a dalších odborných zdrojů.
Beru na~vědomí, že se na moji práci vztahují práva a povinnosti vyplývající
ze zákona č. 121/2000 Sb., autorského zákona v~platném znění, zejména skutečnost,
že Univerzita Karlova má právo na~uzavření licenční smlouvy o~užití této
práce jako školního díla podle §60 odst. 1 autorského zákona.

\vspace{10mm}

\hbox{\hbox to 0.5\hsize{%
V \hbox to 6em{\dotfill} dne \hbox to 6em{\dotfill}
\hss}\hbox to 0.5\hsize{\dotfill\quad}}
\smallskip
\hbox{\hbox to 0.5\hsize{}\hbox to 0.5\hsize{\hfil Podpis autora\hfil}}

\vspace{20mm}
\newpage

%%% Poděkování

\openright

\noindent
\Dedication

\newpage

%%% Povinná informační strana práce

\openright
{\InfoPageFont

\vtop to 0.5\vsize{
\setlength\parindent{0mm}
\setlength\parskip{5mm}

Název práce:
\ThesisTitle

Autor:
\ThesisAuthor

\DeptType:
\Department

Vedoucí \ThesisTypeGenitive{} práce:
\Supervisor, \SupervisorsDepartment

Abstrakt:
\Abstract

Klíčová slova:
{\def\sep{\unskip, }\ThesisKeywords}

\vfil
}

\vtop to 0.49\vsize{
\setlength\parindent{0mm}
\setlength\parskip{5mm}

Title:
\ThesisTitleEN

Author:
\ThesisAuthor

\DeptTypeEN:
\DepartmentEN

Supervisor:
\Supervisor, \SupervisorsDepartmentEN

Abstract:
\AbstractEN

Keywords:
{\def\sep{\unskip, }\ThesisKeywordsEN}

\vfil
}

}

\newpage

%%% Další stránky budeme číslovat
\pagestyle{plain}


%%% Strana s automaticky generovaným obsahem práce

\tableofcontents

%%% Jednotlivé kapitoly práce jsou pro přehlednost uloženy v samostatných souborech
\chapter*{Úvod}
\addcontentsline{toc}{chapter}{Úvod}


\chapter{Základní hra}

\chapter{Algoritmy řešící [n,k]-Mastermind}

\section{Kategorie algoritmů}
% přidat definici algoritmu, který řeší mastermind - vytvoří posloupnost kódů, která končí tajným kódem s nejlepším ohodnocením. 
Algoritmy řešící hru Mastermind lze rozdělit do dvou kategorií, \textbf{deterministické} a \textbf{nedeterministické}. Deterministický algoritmus pro dva shodné stavy volí další krok vždy jednoznačně. Nemůže se tedy stát, že by deterministický algoritmus pro dva stejné vstupy vrátil odlišné výsledky. Nedeterministické algoritmy naproti tomu mohou mít v každém stavu na výběr z více následujících kroků. Volit mohou například podle nějakého pravděpodobnostního rozdělení. Není tedy zaručeno, že pro stejný vstup vrátí vždy identický výsledek.

Hlavní předností deterministických algoritmů je záruka výsledku. Na druhou stranu ale tyto typy algoritmů mohou mít velikou časovou složitost, protože jsou často založeny na procházení všech možností pokračování. To se snaží řešit nedeterministické algoritmy, které hledají rovnováhu mezi časovou složitostí a dobrými výsledky. 

V této práci budeme zkoumat pouze deterministické algoritmy. 

% šlo by zmínit výhody a nevýhody (ne)deterministických algoritmů
% Popsat algoritmy, které hrají pouze kandidáty????


\section{Deterministické algoritmy}
% mozna napsat pozdeji ... Všechny deterministické algoritmy, které budeme srovnávat ...

\subsubsection{stav hry}
V každém kroku algoritmů vycházíme z určitého stavu hry, pro který daná metoda vybírá další pokus. Stav je posloupnost pokusů s ohodnocením a jsou to přesně ty informace, které má hráč [n,k]-Mastermindu k dispozici.

\begin{definice}[Stav]\label{stav}
   Nechť $u_i, i\in \{1,2,\dots m\}$ jsou kódy a $r_i = (b_i,w_i)$ jsou ohodnocení. Stav definujeme jako posloupnost $((u_1, r_1), \dots, (u_m, r_m))$. Počáteční stav definujeme jako prázdnou posloupnost.
\end{definice}

%Pro každý stav lze nalézt množinu kódů, které mají s pokusy ve stavu stejné ohodnocení. 

Pro každý stav lze nalézt množinu kódů, které by podle dostupných informací mohly být tajným kódem. Jde o ty kódy, které mají s kódy ve stavu stejné ohodnocení jako příslušná ohodnocení ve stavu. 

\begin{definice}[kandidát]\label{kandidat}
  Nechť $\left((u_1, r_1), (u_2,r_2), \dots, (u_j,r_j)\right), u_i \in H_{n,k}, r_i \in \N _0 \times \N _0$ je stav. Řekneme, že kód $u \in H_{n,k}$ je kandidát, pokud $s(u,u_i) = r_i \hspace{5px} \forall i \in \{1, \dots j\}$. Pro počáteční stav definujeme množinu kandidátů jako celý prostor kódů $H_{n,k}$.
\end{definice}


\begin{definice}[Prázdný a konečný stav]\label{kandidat}
  Nechť $P = \left((u_1, r_1), (u_2,r_2), \dots, (u_j,r_j)\right), u_i \in H_{n,k}, r_i \in \N _0 \times \N _0$ je stav. Řekneme, že je prázdný, pokud je množina kandidátů tohoto stavu prázdná. 

  Řekneme, že stav $P$ je konečný, pokud je neprázdný a poslední člen je roven $(u,(n,0))$ pro nějaký kód $u$.
\end{definice}

Prázdný stav nemůže odpovídat průběhu hry Mastermind, protože v tu chvíli žádný kód nemůže být tajným kódem. V reálné hře by ho šlo dosáhnout pouze při chybě v ohodnocení pokusu. Konečný stav odpovídá konci hry, kdy se pokus shoduje s tajným kódem. Množina jeho kandidátů v tu chvíli obsahuje právě tajný kód. 

Průběh hry popíšeme jako kořenový strom kde vrcholy odpovídají stavům hry. Hrany znázorňují pokusy s ohodnocením. 

\begin{definice}[Strom \text{[n,k]-Mastermindu}]
  Strom \text{[n,k]-Mastermindu} definujeme jako orientovaný graf na množině stavů. Kořenem je počáteční stav. Ze stavu $A = \left((u_1, r_1), (u_2,r_2), \dots, (u_j,r_j)\right)$ do stavu $B = \left((w_1, s_1), (w_2,s_2), \dots, (w_l,s_l)\right)$ vede hrana právě tehdy, když $l = j+1$ a $\forall i \in \{1,2,\dots, j\} \colon (u_i, r_i) = (w_i, s_i)$. Tuto hranu značíme $(w_l, s_l)$. 
\end{definice}

Stavy následující po stavu $A$ nazýváme potomky $A$.

\begin{definice}[Potomek]\label{potomek}
  % Uvažujme [n,k]-Mastermind s nějakým tajným kódem $v\in H_{n,k}$. 
  % Nechť $S$ je množina všech ohodnocení v $H_{n,k}$. Nechť $A = \left((u_1, r_1), \dots, (u_j,r_j)\right)$ je stav, $u \in H_{n,k}$ je kód a $r \in S$ ohodnocení. Potomka stavu $A$ vzhledem ke kódu $u$ a ohodnocení $r$ definujeme jako stav $A_{u,r} = (\left((u_1, r_1), \dots, (u_j,r_j) ,(u,r)\right)$. 
  
  Potomek stavu $A$ vzhledem ke kódu $u$ a ohodnocení $r$ je stav, na který se ze stavu $A$ dostaneme hranou $(u,r)$. 
  
  Nechť $K$ je množina kódů, potom potomka $K$ vzhledem ke kódu $u$ a ohodnocení $r$ definujeme jako množinu
  \[K_{u,r} = \{w \in K \mid s(u,w) = r\}.\] 
\end{definice}

\begin{pozn}
    Potomkem $K$, vzhledem ke kódu $u$ nazýváme množinu $K_{u,r}$ pro nějaké $r \in S$. Potomkem $K$ nazýváme množinu $K_{u,r}$ pro nějaké $u\in H_{n,k}$ a $r \in S$. 
\end{pozn}



\begin{pozn}
    Používáme stejnou definici pro dva různé objekty. Jejich význam je ale velice podobný. Označíme-li $K$ množina kandidátů stavu $A$, tak množina kandidátů stavu $A_{u,r}$ je právě právě $K_{u,r}$. Množina $K_{u,r}$ ale může odpovídat potomkům stavu vzhledem k různým kódům. 
    %Platí, že množina kandidátů potomka stavu je potomek množiny. Potomek množiny ale může odpovídat více stavům. 
\end{pozn}

\begin{pozn}
    Všimněme si, že $K_{u,r} \subset K \hspace{3px}\forall u\in H_{n,k} \hspace{3px}\forall r\in S$. Navíc pro každou množinu $K \subset H_{n,k}$ a kód $u \in H_{n,k}$ platí
    \[K = \bigcup_{r\in S} K_{u,r}\]
    protože pro každý kód $w \in K$ platí, že $w \in K_{u, s(u,w)}$ z definice potomka.
    %, tedy že $w$ náleží do potomka $K$ vzhledem ke kódu $u$ a ohodnocení $s(u,w)$. 
\end{pozn}

%\begin{pozn} Platí následující pozorování.\[K_{u,r} = \{w \in K \mid s(u,w) = r\}.\] \end{pozn}
% Šlo by definovat funkci, která stavu přiřadí množinu handidátů. Valuace a strategie by ale šly zachovat pro množinu kandidátů. 
Obrázek \ref{fig22prvnitah} zobrazuje počáteční stav a všechny jeho potomky. Jsou to všechny stavy, do kterých se lze dostat po prvním pokusu. 


\begin{figure}[h!]
    \centering
    \begin{tikzpicture}
    \node (1) at (0,0) {\{11,12,21,22\}};
    
    \node (2) at (-3,-3) {\{22\}};
    \draw[->,blue] (1) -- (2) node[pos=0.5, below, sloped] {$11,(0,0)$};
    \node (3) at (0,-3) {\{12, 21\}};
    \draw[->,blue] (1) -- (3) node[pos=0.5, below, sloped] {$11,(1,0)$};
    \node (4) at (3,-3) {\{11\}};
    \draw[->,blue] (1) -- (4) node[pos=0.5, below, sloped] {$11,(2,0)$};
    
    \node (5) at (6,-2) {\{21\}};
    \draw[->,black!60!green] (1) -- (5) node[pos=0.5, above, sloped] {$12,(0,2)$};
    \node (6) at (6,0) {\{11, 22\}};
    \draw[->,black!60!green] (1) -- (6) node[pos=0.5, above, sloped] {$12,(1,0)$};
    \node (7) at (6,2) {\{12\}};
    \draw[->,black!60!green] (1) -- (7) node[pos=0.5, above, sloped] {$12,(2,0)$};

    \node (8) at (-3,3) {\{12\}};
    \draw[->] (1) -- (8) node[pos=0.5, above, sloped] {$21,(0,2)$};
    \node (9) at (0,3) {\{11, 22\}};
    \draw[->] (1) -- (9) node[pos=0.5, above, sloped] {$21,(1,0)$};
    \node (10) at (3,3) {\{21\}};
    \draw[->] (1) -- (10) node[pos=0.5, above, sloped] {$21,(2,0)$};

    \node (11) at (-6,-2) {\{11\}};
    \draw[->,red] (1) -- (11) node[pos=0.5, above, sloped] {$22,(0,0)$};
    \node (12) at (-6,0) {\{12, 21\}};
    \draw[->,red] (1) -- (12) node[pos=0.5, above, sloped] {$22,(1,0)$};
    \node (13) at (-6,2) {\{22\}};
    \draw[->,red] (1) -- (13) node[pos=0.5, above, sloped] {$22,(2,0)$};


    \end{tikzpicture}
    \caption{Graf všech potomků $H_{2,2}$}
    \label{fig22prvnitah}
\end{figure}

Průběh [n,k]-Mastermindu s nějakým tajným kódem $v$ lze sledovat ve stromu [n,k]-Mastermindu. Nechť $A$ je stav a hráč zvolí další pokus $u$. Touto volbou hráč vybral další vrstvu stavů hry, kterou jsou právě potomci stavu $A$ vzhledem ke kódu $u$. Zadavatel ohodnotí kód $u$ vzhledem ke kódu $v$ ohodnocením $r$. Ohodnocením určí, který stav z potomků $A$ vzhledem ke kódu $u$ bude následovat. Hráč tedy může vybrat následující pokus ten, který bude mít nejvhodnější vrstvu potomků vzhledem ke stavu $A$. Z těchto stavů ale neví, do kterého se hra dostane, protože nezná ohodnocení s tajným kódem. 

% Kdybychom si zapisovali uvedený strom pro všechny tajné kódy $v \in H_{n,k}$, tedy bychom se mohli do všech následujících neprázdných stavů, pomocí různých tajných kódů.  sestrojili bychom takzvaný strom algoritmu. 

Předpokládejme, že máme nějaký algoritmus výběru dalšího tahu a pro každý dosažitelný stav [n,k]-Mastermindu s libovolným kódem $v$ známe následující zahraný kód. Potom lze sestavit takzvaný strom algoritmu, který zobrazuje pro každý stav potomky vzhledem ke zvolenému kódu pro tento stav. Další pokus má smysl hledat pouze pro neprázdné stavy, které nejsou konečné. Konečné stavy totiž odpovídají skončené hře. 

% Zajímají nás všichni potomci až do chvíle, kdy dospějeme do konečného stavu, který odpovídá dohrané hře. 

% Zadavatel ohodnocením kódu ... Toto ohodnocení následně určí zadavatel. ...Následně zahraný kód s ohodnocením jednoznačně definují další stav. ...který udělením ohodnocení určí, který z potomků odpovídá ohodnocení vzhledem k tajnému kódu. ...Zadavatel následně určí, který z potomků odpovídá ohodnocení vzhledem k tajnému kódu. 

% stavu, kterým je potomek $A$ vzhledem ke kódu $u$ a ohodnocení $r$. 


% Nyní definujeme strom algoritmu, který odpovídá průběhu hry při pevně zvolené metodě výběru dalších pokusů. Stavy se větví podle různých ohodnocení pokusů. 

\begin{definice}[Strom algoritmu \text{[n,k]-Mastermindu}]
  Strom algoritmu definujeme jako souvislý podgraf stromu [n,k]-Mastermindu, který splňuje následující podmínky.
  \begin{enumerate}
      \item Z každého vrcholu vedou hrany pouze k neprázdným potomkům vzhledem k jednomu určitému kódu pro každý stav. 
      % \item Obsahuje $k^n$ listů, které jsou konečnými stavy.
      \item Pro každý kód $u \in H_{n,k}$ existuje právě jeden list, který je konečným stavem s posledním členem $(u, (n,0))$. 
  \end{enumerate}
\end{definice}

% Nechť máme nějaký algoritmus, který pro každý stav vybere deterministicky nějaký další pokus. Strom algoritmu lze sestrojit rekurzivně. Pro každý stav algoritmus nalezne další pokus. Do stromu zakreslíme všechny potomky stavu vzhledem k tomuto pokusu.

\begin{pozn}
    % Podmínka souvislosti grafu značí, že zapisujeme stavy, do který se následně lze dostat cestou z počátečního stavu. První podmínka zaručuje, že pro každý stav lze zahrát maximálně jeden kód. 
    
    Nechť $A$ je stav. Platí, že pro nějaký kód $u$ jsou množiny kandidátů potomků $A$ vzhledem ke kódu $u$ disjunktní a sjednocením dávají množinu kandidátů stavu $A$. Všechny stavy, do kterých se dostaneme cestou z nějakého stavu $A$ za dodržení podmínky $1$ tedy mají kandidáty z množiny kandidátů stavu $A$. 

    Díky konstrukci stromu [n,k]-Mastermindu strom algoritmu neobsahuje žádné cykly a díky souvislosti a orientaci hran obsahuje nějaký kořen. Pokud by kořenem nebyl počáteční stav, množina kandidátů kořene by nebyla $H_{n,k}$. Proto by nešlo najít cestu do konečného stavu s nějakým kódem $u$, a tedy by nebyla splněna podmínka $2$. 

    Nechť nyní $v\in H_{n,k}$ je nějaký tajný kód a máme nějaký algoritmus popsaný stromem algoritmu. Potom průběh hry je popsán právě cestou z počátečního stavu do konečného stavu s posledním prvkem $(v,(n,0))$. Příklad stromu algoritmu pro [2,2]-Mastermind ukazuje obrázek \ref{fig22minmax}. 
    
    % Použitím obou podmínek tedy dospějeme k závěru, že strom algoritmu musí obsahovat počáteční stav. 

    % Strom obsahuje nějaký kořen, protože je souvislý a ve stromu [n,k]-Mastermindu neexistují cykly. 
    
    %Druhou podmínkou zaručíme, že existuje průběh hry od počátečního stavu, který 
\end{pozn}


%\begin{pozn}
 %   Podstrom algoritmu přesně popisuje zahrané tahy ve všech možných stavech hry. Z každého vrcholu, který není listem vede hrana do všech neprázdných potomků vzhledem k danému kódu. 
%\end{pozn}



%\begin{pozn}
    %Stavy budeme také označovat množinou kandidátů pro tento stav. Množina kandidátů neurčuje jednoznačně stav (například prohozením prvků stavu bychom došli do stejné množiny kandidátů). Naopak stav jednoznačně určuje množinu kandidátů. 
%\end{pozn}

\begin{figure}[h!]
    \centering
    \begin{tikzpicture}
    \node (1) at (0,0) {\{11,12,21,22\}};
    \node (2) at (-4,-2) {$((11,(0,0)))$};
    \draw[->] (1) -- (2) node[pos=0.5, left] {$11,(0,0)$};
    \node (3) at (0,-2) {$(11,(1,0))$};
    \draw[->] (1) -- (3) node[pos=0.5, right] {$11,(1,0)$};
    \node (4) at (4,-2) {$(11,(2,0))$};
    \draw[->, thick, blue] (1) -- (4) node[pos=0.5, right] {$11,(2,0)$};
    \node (5) at (-1,-4) {$((11,(1,0)), (12,(0,2)))$};
    \draw[->] (3) -- (5) node[pos=0.5, left] {$12,(0,2)$};
    \node (6) at (4,-4) {$((11,(1,0)), (12,(2,0)))$};
    \draw[->, thick, blue] (3) -- (6) node[pos=0.5, right] {$12,(2,0)$};
    \node (7) at (-6,-4) {$((11,(1,0)), (22,(2,0)))$};
    \draw[->, thick, blue] (2) -- (7) node[pos=0.5, right] {$22,(2,0)$};
    \node (8) at (-1,-6) {$((11,(1,0)), (12,(0,2)), (21,(2,0)))$};
    \draw[->, thick, blue] (5) -- (8) node[pos=0.5, left] {$21,(2,0)$};
    \end{tikzpicture}
    \caption{Strom algoritmu pro [2,2]-Mastermind}
\label{fig22minmax}
\end{figure}


% možná to není potřeba
%\begin{lemma}[Jednoznačnost hrany]
 %   Nechť $(V, E)$ je graf a $K \in V$ je vrchol. Potom z $K$ vede právě jedna hrana 
%\end{lemma}



% tato definice možná ani není potřeba, protože s ní dále nepracuji
%\begin{definice}[Rozdělení]
 %   Pro $K \subset H_{n,k}$ definujeme rozdělení vzhledem ke kódu $u$ jako množinu potomků $K$, vzhledem ke kódu $u$. $R_{K,u} = \{K_{u,r} \mid r \in S\}$. 
%\end{definice}


\subsubsection{Obecný algoritmus}
Nyní přistoupíme k popisu skupiny deterministických algoritmů řešící hru [n,k]-Mastermind. Využijeme k tomu dvě funkce, valuaci a strategii.
% U této definice si nejsem jistý, jak to smyslupně definovat. Vazba "která lz e vyjádřit pomocí potomků" se mi moc nelíbí - definovat valuaci a až potom jednokrokové valuace
\begin{definice}[Valuace]
    Nechť $G = (V,E)$ je graf [n,k]-Mastermindu a $K \in V$ je vrchol. Potom valuaci vzhledem k množině $K$ definujeme jako funkci $f_K(u) \colon H_{n,k} \to \mathbb{R}$.
\end{definice}
Nejprve si musíme uvědomit, že pro dva různé stavy, které ale mají shodnou množinu kandidátů je problém nalezení dalšího kódu ekvivalentní. Předpokládáme totiž, že všechny kódy mohou být tajným kódem se stejnou pravděpodobností. Stejné množiny kandidátů pro dva různé stavy lze například dosáhnout prohozením prvků stavu. 

Valuace bude v algoritmech sloužit pro vyčíslení vhodnosti kódu $u$ jako dalšího pokusu pro stav vyjádřený množinou kandidátů $K\subset H_{n,k}$. Tato hodnota bude obvykle záviset na potomcích $K$ vzhledem ke kódu $u$. Potomci $K$ vzhledem ke kódu $u$ totiž určují možné následující stavy v případě, kdy zvolíme kód $u$ jako další pokus. Valuace nám umožňuje ohodnotit kód podle toho, jak dobré jsou jeho následující stavy. 

\begin{prikl}\label{prjednokrokfce}
    Funkce $f_K$ je příklad valuace, která kódu $u$ přiřadí očekávaný počet kódů potomka $K$ vzhledem ke kódu $u$.
    % jednokrokov8 funkce je taková, která kódu $u$ přiřadí hodnotu, která lze vyjádřit pomocí potomků $K$ vzhledem ke kódu $u$.
    \begin{align*}
        f_K \colon H_{n,k} &\to \mathbb{R} \\
        u &\mapsto \sum_{r\in S}\frac{|K_{u,r}|}{|K|}|K_{u,r}| 
    \end{align*}
    Pro $K = H_{2,2}$ z obrázku \ref{fig22prvnitah} má $f_{K}$ konstantní hodnotu $\frac{3}{2}$.
\end{prikl}
% množinu kandidátů na tajný kód
Obecně valuace nemusí být omezená na následující vrstvu. Mohla by odkazovat na očekávaný počet kódů vrcholu o dvě či více vrstev dál. 

\begin{definice}[Strategie]
    Nechť $\mathcal{F} = \{f_K\colon H_{n,k} \to \mathbb{R} \mid K \subset H_{n,k}\}$ je prostor valuací. Potom strategii definujeme jako zobrazení $F \colon \mathcal{F} \to \mathbb{R}$. 
\end{definice}
% Strategie určitým způsobem slouží k výběru následujícího kódu. 
%Zjednodušeně řečeno slouží k určení, 
Strategie v této práci bude určovat, zda chceme vybírat kódy s vyšší nebo nižší hodnotou valuace.
%Použití bude popsáno níže v popisu algoritmu \ref{alg-default}. 

\begin{prikl}\label{prstrategie}
    Funkce $F$ je strategie, která funkci $f_K \in \mathcal{F} = \{f_K\}$ přiřadí její maximální hodnotu na $H_{n,k}$.
    \[F(f_K) =  \max_{u\in H_{n,k}}\{f_K(u)\}\]
    Strategie je dobře definovaná, protože maximum na konečné množině reálných čísel má vždy jednoznačnou hodnotu.
\end{prikl}

Deterministické algoritmy popisujeme jako algoritmy, které pro následující pokus vybírají kód podle zvolené valuace a strategie. Postup je znázorněn v algoritmu \ref{alg-default}.

% \subsection{Varianty algoritmů}





%Nechť \[P = \left((u_1, r_1), (u_2,r_2), \dots, (u_j,r_j)\right), u_i \in H_{n,k}, r_i \in \N _0 \times \N _0\] jsou tahy s ohodnocením a $K$ je konec cesty $P$ z vrcholu $H_{n,k}$. 

%Potom jednokrokové strategie volí další kód ten, který minimalizuje respektive maximalizuje funkci $f_K$. Pokud je těchto kódů více, zvolí kód, který je zároveň kandidát cesty $P$. Pokud výběr stále není jednoznačný, algoritmus zvolí lexikograficky nejnižší kód. 

\begin{algorithm}[h!]
\begin{algorithmic}[1]  % [1] způsobí, že číslujeme kroky algoritmu
\Function{Solve$[n, k, \mathcal{F}, F]$}{$v$}
    \State $K \gets H_{n,k}$ 
    \State $j \gets 0$
    \State $r_0 \gets (0,0)$
    \While {$r_j \neq (n,0)$} \hfill \mbox{(dokud hra není dohrána)}
        \State $j \gets j + 1$ 
	\State $U \gets \{u \in H_{n,k} \mid f_K(u) = F(f_K)\}$
        \If{$U \cap K \neq \emptyset$}
            \State $u_j \gets$ lexikograficky nejmenší $u \in U \cap K$
	\Else
		\State $u_j \gets$ lexikograficky nejmenší $u \in U$
	\EndIf
        \State $r_j \gets s(u_j, v)$ \hfill \mbox{(ohodnocení pokusu)}
        \State $K \gets K_{u_j,r_j}$
    \EndWhile
    \State \Return $u_j, j$ \hfill \mbox{($u_j$ je tajný kód, $j$ je počet pokusů)}
\EndFunction
\end{algorithmic}
\caption{Algoritmus řešící [n,k]-Mastermind}
\label{alg-default}
\end{algorithm}

Písmenem $K$ značíme množinu kandidátů na tajný kód pro aktuální stav hry (cestu pokusů s ohodnocením). Začínáme s $K = H_{n,k}$. Aktuální číslo pokusu značí $j$, $u_j$ je kód, který algoritmus zvolí jako $j$-tý pokus a $r_j$ je ohodnocení $u_j$ vzhledem k tajnému kódu $v$. Množina $U$ značí kódy, ve kterých valuace $f_K$ nabývá určitou optimální hodnotu danou strategií $F$. Z této množiny algoritmus vybírá další pokus $u_j$. Ve chvíli, kdy existuje kandidát v množině $U$ ($U\cap K \neq \emptyset$), algoritmus zvolí ten, který je lexikograficky nejmenší. Jinak algoritmus volí lexikograficky nejmenší kód z $U$. 

Aby byla strategie $F$ použitelná v algoritmu, musí platit následující podmínka. 
\[\forall K \in V \hspace{5px} \exists u \in H_{n,k} \colon F(f_K) = f_K(u)\]
Pokud by tato podmínka nebyla splněna, mohlo by se stát, že množina $U$ bude prázdná a algoritmus nedokáže zvolit další pokus.

% šlo by napsat pro dvojice funkce f, funkcionál F
%\begin{veta}[Správnost algoritmu \ref{alg-default}]
 %   Algoritmus \ref{alg-default}
%\end{veta}

\begin{tvrz}[Počet ohodnocení]
Nechť $n\in \mathbb{N}, k\in \mathbb{N}, k \geq 3$. Potom počet všech možných ohodnocení kódů v $H_{n,k}$ je $\frac{n^2 + 3n}{2}$.
\end{tvrz}
\begin{dukaz}
Pro $b \leq n,\hspace{3px} b \neq n-1$ černých kolíčků existuje $n-b+1$ možností na počet bílých kolíčků podle tvrzení \ref{tvrzohodnoceni}. Pro $b = n-1$ existuje pouze jedno ohodnocení. Součtem přes počet černých kolíčků dostáváme následující počet možností ohodnocení.
    \[|S| = (n+1) + n + (n-1)  + \dots + 3 + 1 + 1 \]
Vzorcem pro součet aritmetické řady se dobereme k výsledku.
    \[|S| = \frac{(n+1)(n+2)}{2} - 1 \]
    \[|S| = \frac{n^2 + 3n}{2}\]
\end{dukaz}

\begin{tvrz}[Časová složitost algoritmu \ref{alg-default}]
    Uvažujeme případ, kdy $F$ je minimizér/maximizér.
    Algoritmus \ref{alg-default} má časovou složitost
    % \[O( k^n O(f_K))\]
    \[O \left( \sum_{j = 1}^z k^n \cdot m_j \cdot n^2 \cdot \frac{n^2 + 3n}{2}\right)\]
    
\end{tvrz}
\begin{dukaz}
    $k^n$ je počet všech kódů, které musím projít, abych našel ten nejlepší, $m_j$ je maximální počet kandidátů v j-té iteraci, $z$ je maximální počet iterací algoritmu. $\frac{n^2 + 3n}{2}$ je počet všech ohodnocení, $n^2$ je časová složitost ohodnocení.

    Potřebuji ještě limit na počet iterací $z$, horní odhad na $m_j$ a zjistit složitost nalezení průniku $U$ a $K$.
\end{dukaz}



\chapter{Jednokrokové algoritmy}
Jednokrokové algoritmy jsou ty, jejichž valuace lze vyjádřit pouze pomocí následující vrstvy potomků. Tedy ohodnocuje kód v závislosti na rozdělení potomků vzhledem k tomuto kódu. V této kapitole porovnáme tři jednokrokové algoritmy, které se liší pouze podle zvoleného prostoru valuací a strategie. Jsou to algoritmy Min-max, Max entropy a Most parts. 

%definujeme následovně. Nechť $G = (V,E)$ je graf [n,k]-Mastermindu. Nechť $\mathcal{F} = \{f_K \mid K \in V \} $ je prostor jednokrokových funkcí, kde $f_K \colon H_{n,k} \to \mathbb{R}$ jsou jednokrokové funkce. Nechť $F$ je funkcionál na $\mathcal{F}$. Potom jednokrokový algoritmus s prostorem funkcí $\mathcal{F}$ a funkcionálem $F$ je definovaný v algoritmu \ref{alg-default}. 



% Nějak popsat to, jak odhadujeme výsledný počet tahů podle tabulky rozdělení - max, entropie, počet částí, očekávaná velikost části v tabulce rozdělení - Irwing, nezmínil jsem.

\section{Min-max}
Prvním algoritmem je metoda, kterou navrhl D. Knuth \cite{donald_e__knuth_1977}. Jeho algoritmus pro libovolný stav $A$ s množinou kandidátů $K$ vybírá jako další pokus ten kód, který minimalizuje maximální počet kódů v potomcích $K$ vzhledem k nějakému kódu. 
% Odpovídá algoritmu \ref{alg-default} se vstupním prostorem valuací $\mathcal{M}$ a strategií $L$. 
%Tedy minimalizuje maximální počet kódů v nějakém stavu, do kterého se hra může dostat po zahrání dalšího kódu.
\begin{definice}[Maximum potomka]
    Nechť $H_{n,k}$ je prostor kódů. Pro $K \subset H_{n,k}$ definujeme funkci maximum potomka jako
    \begin{align*}
        m_K \colon H_{n,k} &\to \mathbb{R} \\
        u &\mapsto \max_{r\in S} |K_{u,r}|.
    \end{align*}
    Prostor všech funkcí $m_K$ označíme $\mathcal{M}$. 
\end{definice}

\begin{definice}[Strategie minimum]
    Nechť $H_{n,k}$ je prostor kódů a $\mathcal{F} = \{f_K\colon H_{n,k} \to \mathbb{R} \mid K \subset H_{n,k}\}$ je prostor valuací. Potom na $\mathcal{F}$ definujeme strategii minimum $L$ jako
    \begin{align*}
        L \colon \mathcal{F} &\to \mathbb{R} \\
        f &\mapsto \min_{u\in H_{n,k}} f(u)
    \end{align*}
\end{definice}

\begin{definice}[Min-max algoritmus]
    \emph{Min-max algoritmus} definujeme jako deterministický algoritmus s prostorem valuací $\mathcal{M}$ a strategií $L$. Tedy je to funkce, která tajnému kódu $v$ přiřadí $\textsc{Solve}[n, k, \mathcal{M}, L](v)$, posloupnost zahraných kódů s ohodnoceními.
\end{definice}


% \begin{definice}[Min-max algoritmus]
%     Nechť $n\in \N, k\in \N$, $\mathcal{M}$ je prostor valuací maxima potomků a $M$ je strategie maximum. Min-max algoritmus je definovaný jako funkce 
%     \begin{align*}
%         \textsc{Solve}[n, k, \mathcal{M}, L] \colon H_{n,k} &\to \mathbb{N} \\
%         v & \mapsto \textsc{Solve}[n, k, \mathcal{M}, L](v)
%     \end{align*}
%     podle předpisu algoritmu \ref{alg-default}.
% \end{definice}

Algoritmus Min-max pro libovolný stav $A$ s množinou kandidátů $K$ vybírá další pokus z množiny
\[U = \{u \in H_{n,k} \mid \max_{r\in S} |K_{u,r}| = \min_{w \in H_{n,k}} \max_{r\in S} |K_{w,r}| \}.\]

\begin{tvrz}[Správnost Min-max algoritmu]
    Min-max algoritmus řeší hru [n,k]-Mastermind.
\end{tvrz}
\begin{dukaz}
    Jednoznačnost řešení plyne z předpisu valuace a strategie a z jednoznačnosti ohodnocení. Dále dokážeme, že v každém kroku algoritmu se zmenší velikost množiny kandidátů aktuálního stavu, a tedy algoritmus skončí. Nechť $K$ je množina kandidátů obsahující alespoň dva kódy. Nechť $u \in K$. Potom $|K_{u,(n,0)}| = 1$, a tedy 
    \[|\bigsqcup_{r\in S, r \neq (n,0)} K_{u,r} | = |K| - 1.\]
    Díky tomu platí následující nerovnost
    \[\min_{w \in H_{n,k}} \max_{r\in S} |K_{w,r}| \leq |K|-1.\]
    V každém kroku se tedy velikost množiny kandidátů aktuálního stavu snižuje. Pokud $|K| = 1$, algoritmus zvolí jako další pokus jedinný zbývající kód z množiny kandidátů, protože pro všechny kódy platí, že
    $\max_{r\in S} |K_{w,r}| = 1$.
    Zároveň tajný kód náleží v každém kroku do množiny kandidátů, a tedy jej algoritmus nalezne. 
\end{dukaz}



Pro [2,2]-Mastermind sestrojíme strom Min-max algoritmu a nastíníme na něm chod Knuthova algoritmu.


\subsubsection{[2,2]-Mastermind}
Prostor kódů [2,2]-Mastermindu je množina $H_{2,2} = \{11,12,21,22\}$. Množina všech ohodnocení je podle tvrzení \ref{tvrzohodnoceni2} rovna $S = \{(0,0), (0,2), (1,0), (2,0)\}$. Na obrázku \ref{fig22prvnitahmnoziny} jsou zobrazeni potomci $H_{2,2}$ vzhledem ke všem kódům a ohodnocením. Pro přehlednost jsou vrcholy uspořádány podle zahraných kódů. Maximální počty prvků ve vrcholech potomků $H_{2,2}$ (hodnoty $m_{H_{2,2}}$) jsou následující. 
\[m_{H_{2,2}} (11) = m_{H_{2,2}} (12) = m_{H_{2,2}} (21) = m_{H_{2,2}} (22) = 2\]
Knuthův algoritmus tedy volí jako první pokus kód $11$. 

\begin{figure}[h!]
    \centering
    \begin{tikzpicture}
    \node (1) at (0,0) {\{11,12,21,22\}};
    \node (2) at (-3,-3) {\{22\}};
    \draw[->,blue] (1) -- (2) node[pos=0.5, below, sloped] {\footnotesize$11,(0,0)$};
    \node (3) at (0,-3) {\{12, 21\}};
    \draw[->,blue] (1) -- (3) node[pos=0.5, below, sloped] {\footnotesize$11,(1,0)$};
    \node (4) at (3,-3) {\{11\}};
    \draw[->,blue] (1) -- (4) node[pos=0.5, below, sloped] {\footnotesize$11,(2,0)$};
    \node (5) at (6,-2) {\{21\}};
    \draw[->,black!60!green] (1) -- (5) node[pos=0.5, above, sloped] {\footnotesize$12,(0,2)$};
    \node (6) at (6,0) {\{11, 22\}};
    \draw[->,black!60!green] (1) -- (6) node[pos=0.5, above, sloped] {\footnotesize$12,(1,0)$};
    \node (7) at (6,2) {\{12\}};
    \draw[->,black!60!green] (1) -- (7) node[pos=0.5, above, sloped] {\footnotesize$12,(2,0)$};
    \node (8) at (-3,3) {\{12\}};
    \draw[->] (1) -- (8) node[pos=0.5, above, sloped] {\footnotesize$21,(0,2)$};
    \node (9) at (0,3) {\{11, 22\}};
    \draw[->] (1) -- (9) node[pos=0.5, above, sloped] {\footnotesize$21,(1,0)$};
    \node (10) at (3,3) {\{21\}};
    \draw[->] (1) -- (10) node[pos=0.5, above, sloped] {\footnotesize$21,(2,0)$};
    \node (11) at (-6,-2) {\{11\}};
    \draw[->,red] (1) -- (11) node[pos=0.5, above, sloped] {\footnotesize$22,(0,0)$};
    \node (12) at (-6,0) {\{12, 21\}};
    \draw[->,red] (1) -- (12) node[pos=0.5, above, sloped] {\footnotesize$22,(1,0)$};
    \node (13) at (-6,2) {\{22\}};
    \draw[->,red] (1) -- (13) node[pos=0.5, above, sloped] {\footnotesize$22,(2,0)$};
    \end{tikzpicture}
    \caption{Graf všech potomků $H_{2,2}$}
    \label{fig22prvnitahmnoziny}
\end{figure}


Nyní algoritmus volí druhý pokus. Pro $K \neq \emptyset$ a $u \in H_{n,k}$ platí, že $m_K(u) \geq 1$. Proto, když lexikograficky nejnižší kandidát $u$ dosahuje hodnoty $m_K(u) = 1$, algoritmus Min-max jej vždy vybere. Konkrétně to bude platit pro jednoprvkovou množinu kandidátů $\{22\}$. Algoritmus zde jako druhý pokus vybírá jediného zbývajícího kandidáta. 

V případě, že ohodnocení kódu $11$ v prvním pokusu je $(1,0)$, nacházíme se ve stavu $\{12,21\}$. Potomci vrcholu $\{12,21\}$ jsou znázorněni na obrázku \ref{fig22druhytah}. 
\begin{figure}[h!]
    \centering
    \begin{tikzpicture}
    \node (3) at (0,0) {\{12, 21\}};

    \node (5) at (4,0) {\{12, 21\}};
    \draw[->,black!60!green] (3) -- (5) node[pos=0.5, above,sloped] {\footnotesize$22,(1,0)$};
    
    \node (5) at (-2,-2) {\{21\}};
    \draw[->,blue] (3) -- (5) node[pos=0.5, above,sloped] {\footnotesize$12,(0,2)$};
    \node (6) at (2,-2) {\{12\}};
    \draw[->,blue] (3) -- (6) node[pos=0.5, above,sloped] {\footnotesize$12,(2,0)$};

    \node (7) at (-2,2) {\{12\}};
    \draw[->] (3) -- (7) node[pos=0.5, above,sloped] {\footnotesize$21,(0,2)$};
    \node (8) at (2,2) {\{21\}};
    \draw[->] (3) -- (8) node[pos=0.5, above,sloped] {\footnotesize$21,(2,0)$};

    \node (9) at (-4,0) {\{12, 21\}};
    \draw[->,red] (3) -- (9) node[pos=0.5, above,sloped] {\footnotesize$11,(1,0)$};

    \end{tikzpicture}
    \caption{Výběr druhého tahu}
\label{fig22druhytah}
\end{figure}
Hodnoty valuací jednotlivých kódů jsou následující:
\[m_{\{12,21\}} (11) = m_{\{12,21\}} (22) = 2, \hspace{10px} m_{\{12,21\}} (12) = m_{\{12,21\}} (21) = 1.\]
Algoritmus Min-max tedy v tomto stavu vybírá jako druhý pokus kód $12$. 

Na obrázku \ref{fig22minmax} je zobrazen strom Min-max algoritmu. Protože vrchol ve stromu algoritmu je roven obarvení cesty z počátku do tohoto vrcholu, pro znázornění postupu algoritmu vrcholy obarvujeme množinou kandidátů příslušného stavu. Algoritmus Min-max vyhraje [2,2]-Mastermind vždy na maximálně tři tahy. 
\begin{figure}[h!]
    \centering
    \begin{tikzpicture}
    \node (1) at (0,0) {\{11,12,21,22\}};
    \node (2) at (-4,-2) {\{22\}};
    \draw[->] (1) -- (2) node[pos=0.5, above, sloped] {\footnotesize$11,(0,0)$};
    \node (3) at (0,-2) {\{12, 21\}};
    \draw[->] (1) -- (3) node[pos=0.5, right] {\footnotesize$11,(1,0)$};
    \node (4) at (4,-2) {\{11\}};
    \draw[->, thick, blue] (1) -- (4) node[pos=0.5, above, sloped] {\footnotesize$11,(2,0)$};
    \node (5) at (-2,-4) {\{21\}};
    \draw[->] (3) -- (5) node[pos=0.5, above, sloped] {\footnotesize$12,(0,2)$};
    \node (6) at (2,-4) {\{12\}};
    \draw[->, thick, blue] (3) -- (6) node[pos=0.5, above, sloped] {\footnotesize$12,(2,0)$};
    \node (7) at (-6,-4) {\{22\}};
    \draw[->, thick, blue] (2) -- (7) node[pos=0.5, above, sloped] {\footnotesize$22,(2,0)$};
    \node (8) at (-2,-6) {\{21\}};
    \draw[->, thick, blue] (5) -- (8) node[pos=0.5, left] {\footnotesize$21,(2,0)$};
    \end{tikzpicture}
    \caption{Zahrané tahy Min-max algoritmem}
\label{fig22minmax}
\end{figure}

% Přidat pseudokód algoritmu???

%vysledky Knuthova algoritmu:
% zdroj - Kooi
%expected - 4.476
% maximal - 5
% použít svůj kód a ocitovat

\section{Maximální entropie}
Neuwirth \cite{neuwirth} navrhl algoritmus hrající kódy, které maximalizují entropii velikostí množin potomků vzhledem ke kódu. Nejprve definujeme valuaci a strategii použité v tomto algoritmu. Následně vysvětlíme myšlenku této strategie.

\begin{definice}[Entropie potomků]\label{defentropierozdeleni}
    Nechť $H_{n,k}$ je prostor kódů. Pro $K \subset H_{n,k}$ definujeme valuaci entropie potomků jako funkci
    \begin{align*}
        h_K \colon H_{n,k} &\to \mathbb{R} \\
        u &\mapsto \sum_{r\in S, |K_{u,r}| > 0} \frac{|K_{u,r}|}{|K|}\log_2\left( \frac{|K|}{|K_{u,r}|} \right).
    \end{align*}
    Prostor všech valuací $h_K$ označíme $\mathcal{H}$. 
\end{definice}

\begin{definice}[Strategie maximum]\label{defstrategiemaximum}
    Nechť $H_{n,k}$ je prostor kódů a $\mathcal{F} = \{f_K\colon H_{n,k} \to \mathbb{R} \mid K \subset H_{n,k}\}$ je prostor valuací. Potom na $\mathcal{F}$ definujeme funkcionál maximum $M$ následujícím předpisem.
    \begin{align*}
        M \colon \mathcal{F} &\to \mathbb{R} \\
        f &\mapsto \max_{u\in H_{n,k}} f(u)
    \end{align*}
\end{definice}


\begin{definice}[Algoritmus Max entropy]
    \emph{Algoritmus Max entropy} definujeme jako deterministický algoritmus s prostorem valuací $\mathcal{H}$ a strategií $M$. Tedy je to funkce, která tajnému kódu $v$ přiřadí posloupnost zahraných kódů s ohodnoceními $\textsc{Solve}[n, k, \mathcal{H}, M](v)$.
\end{definice}


% \begin{definice}[Algoritmus Max entropy]
%     Nechť $n\in \N, k\in \N$, $\mathcal{H}$ je prostor valuací entropie potomků a $M$ je strategie maximum. Algoritmus Max entropy je definovaný jako funkce 
%     \begin{align*}
%         \textsc{Solve}[n, k, \mathcal{H}, M] \colon H_{n,k} &\to \mathbb{N} \\
%         v & \mapsto \textsc{Solve}[n, k, \mathcal{H}, M](v)
%     \end{align*}
%     podle předpisu algoritmu \ref{alg-default}.
% \end{definice}

\begin{tvrz}[Správnost algoritmu Max entropy]
    Max Entropy algoritmus řeší hru [n,k]-Mastermind.
\end{tvrz}
\begin{dukaz}
    Obdobně jako u Min-max algoritmu dokážeme, že se velikost množiny kandidátů v každém kroku zmenší. Nechť $K$ je množina kandidátů obsahující alespoň dva kódy. 
    Nechť $u\in H_{n,k}$ je kód, pro který platí, že $K_{u,r} = K$ pro nějaké $r\in S_{n,k}$. Potom 
    \[h_K(u) = \frac{|K_{u,r}|}{|K|}\log_2\left( \frac{|K|}{|K_{u,r}|} \right) = 0.\]
    
    Pro libovolný kód $u$ platí, že $\log_2\left( \frac{|K|}{|K_{u,r}|} \right)  > 0 $ pro všechna $r\in S_{n,k}, |K_{u,r}| > 0, K_{u,r} \neq K$. Tedy $h_K(u) = 0$ právě tehdy, když $K_{u,r} = K$ pro nějaké $r\in S_{n,k}$. Pokud pro každou alespoň dvouprvkovou množinu existuje kód, pro který $h_K(u) > 0$, v každém kroku algoritmu Max entropy se počet prvků množina kandidátů ostře zmenší. Pro $u \in K$ platí
    \begin{align*}
        h_K(u) &= \sum_{r\in S, |K_{u,r}| > 0} \frac{|K_{u,r}|}{|K|}\log_2\left( \frac{|K|}{|K_{u,r}|} \right) 
    \geq 
    \frac{|K_{u,(n,0)}|}{|K|}\log_2\left( \frac{|K|}{|K_{u,(n,0)}|} \right)
    = \\
    &=\frac{1}{|K|} \log_2 |K| > 0.
    \end{align*}
    Využili jsme toho, že $|K_{u,(n,0)}| = 1$. V každém kroku se tedy velikost množiny kandidátů aktuálního stavu snižuje. Pokud $|K| = 1$, potom $h_K(u) = 0$ pro všechny kódy $u$. V tomto případě algoritmus Max entropy zvolí jediného zbývajícího kandidáta, kterým je tajný kód. 
    
    % Díky tomu, že tajný kód náleží v každém kroku do množiny kandidátů aktuálního stavu, algoritmus jej nalezne na konečný počet pokusů a algoritmus skončí. 
\end{dukaz}

Níže popíšeme, co si pod pojmem entropie můžeme představit a jaký má význam v [n,k]-Mastermindu.

\begin{definice}[Entropie]\label{defentropie}
  Nechť $A$ je konečná množina, $P \colon A \to [0,1]$ je zobrazení takové, že $\sum_{a \in A} P(a) = 1$. Nechť $supp(P) = \{ a \in A \mid P(a) > 0\}$. Potom entropii zobrazení $P$ definujeme jako 
  \[H(P) = \sum_{a \in supp(P)}P(a)\log_2\left(\frac{1}{P(a)}\right)\]
  Zobrazení $P$ budeme říkat rozdělení, dané hodnotě $P(a)$ pravděpodobnost a množině $[0,1]$ pravděpodobnostní prostor. 
\end{definice}

\begin{pozn}
    Definice \ref{defentropierozdeleni} pro pevný kód $u\in H_{n,k}$ koresponduje s definicí entropie pro množinu všech ohodnocení $A = \{s_1, s_2, \dots, s_m \}$ a zobrazení $P_u\colon A \to [0,1], P_u(s_i) = \frac{|K_{u,r}|}{|K|}$.
\end{pozn}

Intuici za definicí entropie nastíníme v příkladu \ref{prbinarysearch}, který dává určitý smysl pro výraz $\log_2\left(\frac{1}{P(a)}\right)$ a příkladu \ref{prentropierozdeleni}, který ukazuje příklad entropie na konkrétním rozdělení.

\begin{prikl}[Binární vyhledávání]\label{prbinarysearch}
Představme si, že chceme uhodnout tajné číslo z množiny $\{1,2,\dots,8\}$ pokládáním nějakých ano/ne otázek. Předpokládejme, že každé číslo může být tajným číslem se stejnou pravděpodobností $\frac{1}{8}$. Známým postupem, jak tajné číslo najít je takzvaným půlením intervalu (binární vyhledávání). Vždy se zeptáme, zda je číslo větší nebo menší než číslo v polovině. Nechť tajné číslo je $4$. Binární vyhledávání by mělo následující průběh:

Otázka: "Je číslo větší než $4$?". 
Odpověď: ne

Otázka: "Je číslo větší než $2$?". 
Odpověď: ano

Otázka: "Je číslo větší než $3$?". 
Odpověď: ano

Výsledné číslo je $4$. Počet otázek nutných k určení tajného čísla je $3 = \log_2 8$.
\end{prikl}

\begin{pozn}\label{poznotazkynamnozinu}
Obdobně jako v příkladu \ref{prbinarysearch} lze pro jakoukoliv osmiprvkovou množinu $A$ s rovnoměrnou pravděpodobností vytvořit sadu otázek ano/ne a o každém prvku této množiny určit, zda odpovídá odpovědi ano nebo odpovědi ne. Následně bychom mohli pro určení jakéhokoliv prvku této množiny aplikovat obdobný postup jako v příkladu \ref{prbinarysearch}. Pro vhodnou sadu otázek existuje postup, který každou otázkou sníží prostor zbývajících možností na polovinu. Tím pádem pro určení jakéhokoliv prvku množiny $A$ potřebujeme $\log_2(|A|)$ otázek. V případě, kdy velikost množiny není mocnina $2$ bychom mluvili o $\lceil \log_2 |A| \rceil$, kde $\lceil \cdot \rceil$ značí horní celou část. 
\end{pozn}

\begin{prikl}[m odpovědí]\label{prmodpovedi}
    Ve chvíli, kdybychom nebyli omezeni na ano/ne otázky, ale mohli bychom pokládat otázky na které bychom dostali $m$ odpovědí, mohli bychom aplikovat obdobný postup jako v příkladu \ref{prbinarysearch} s binárním vyhledáváním. Tentokrát bychom ale mohli po každé otázce snížil prostor zbývajících možností $m$-krát. Dobrali bychom se k výsledku, že očekávaný počet otázek k určení prvku z množiny $A$ by byl $\log_m |A|$.
\end{prikl}

\begin{prikl}\label{prentropierozdeleni}
Nechť $A = \{ \text{a, b, c, d}\}$ a $P$ je zadaná následovně: 
\[P(a) = \frac{1}{8}, P(b) = \frac{1}{8}, P(c) = \frac{1}{4}, P(d) = \frac{1}{2}\]
Nyní si představme pravděpodobnostní prostor jako čtverec na obrázku \ref{prprostor}, kde pravděpodobnost prvku je velikost plochy. Výraz $\frac{1}{P(a)}$ říká, kolik prvků $A$ s touto pravděpodobností by se "vešlo" do pravděpodobnostního prostoru jako na obrázku. Například dílek c by se "vešel" do pravděpodobnostního prostoru čtyřikrát.
Hodnota $\log_2\left(\frac{1}{P(c)}\right)$ lze interpretovat jako počet otázek ano/ne, který potřebujeme k jednoznačnému určení prvku c v případě, kdyby byl pravděpodobnostní prostor rozdělen rovnoměrně. Tento postup byl popsán 
%v příkladu \ref{prbinarysearch} a 
v poznámce \ref{poznotazkynamnozinu}.
Entropie tohoto rozdělení by v tomto případě byla 
\[H(P) = \frac{1}{8}\log_2 (8) + \frac{1}{8}\log_2 (8) + \frac{1}{4}\log_2 (4) + \frac{1}{2}\log_2 (2) = \frac{3}{8} + \frac{3}{8} + \frac{2}{4} + \frac{1}{2} = \frac{7}{4}.\]
Můžeme ji interpretovat jako očekávaný počet otázek, který potřebujeme k jednoznačnému určení prvku z množiny $A$. 


\begin{figure}
\begin{tikzpicture}
    \draw[black, thick] (0,0) rectangle (4,4);
    \begin{scope}
        \node (1) at (1,0.5) {a - 1/8};
        \node (2) at (1,1.5) {b - 1/8};
        \node (3) at (1,3) {c - 1/4};
        \node (4) at (3,2) {d - 1/2};
    \end{scope}
    \draw[black, thick] (2,0) -- (2,4);
    \draw[black, thick] (0,2) -- (2,2);
    \draw[black, thick] (0,1) -- (2,1);

    
\end{tikzpicture}
\caption{Pravděpodobnostní prostor}
\label{prprostor}
\end{figure}

\end{prikl}

\subsubsection{Entropie a Mastermind}
Nyní vysvětlíme, jakým způsobem aplikujeme výpočet entropie na strategii pro řešení [n,k]-Mastermindu. Otázky v předchozích příkladech nahradíme zahranými pokusy. Jako odpovědi budou sloužit ohodnocení pokusů. Nechť $K$ je aktuální množina kandidátů. Pro $K$ můžeme ve smyslu příkladu \ref{prmodpovedi} definovat nějakou abstraktní hodnotu počtu pokusů potřebných k určení jakéhokoliv prvku  $v \in K$ jako $\log_{|S_{n,k}|}(|K|)$. Na každý pokus můžeme totiž dostat maximálně $|S_{n,k}|$ možných ohodnocení. 

Nechť $u \in H_{n,k}$ je kód. 
% Potom $\{K_{u,r} \mid r\in S_{n,k}\}$ je množina potomků $K$ vzhledem ke kódu $u$ a 
Potom $\log_{|S|}(|K_{u,r}|)$ je očekávaný počet pokusů potřebný k určení libovolného prvku $v \in K_{u,r}$. Pomocí těchto hodnot lze vytvořit odhad na počet pokusů potřebný k určení jakéhokoliv kódu $v \in K$ po zahrání kódu $u$ jako jejich vážený průměr:
\begin{align}\label{rceocekavanypocetpokusu}
    G = \sum_{r\in S_{n,k}, |K_{u,r}| > 0} \frac{|K_{u,r}|}{|K|}\log_{|S_{n,k}|}|K_{u,r}|.
\end{align}



%Pro každé $K_i, i \in \{1, 2, \dots, m\}$ můžeme ve smyslu poznámky \ref{poznotazkynamnozinu} definovat nějakou abstraktní hodnotu počtu otázek ano/ne potřebných k určení jakéhokoliv prvku  $x \in K_i$ jako $\log_2(|K_i|)$. Potom obdobně jako v příkladu \ref{prentropierozdeleni} můžeme určit očekávaný počet otázek jako
%\[G = \sum_{i =1}^n \frac{|K_i|}{|K|}\log_2|K_i| \]
% Tato hodnota bude naším stavebním kamenem pro strategii za užití entropie.
Protože cílem [n,k]-Mastermindu je najít tajný kód v co nejmenším počtu pokusů, chceme hodnotu $G$ minimalizovat. V následující větě dokážeme ekvivalenci problému minimalizace $G$ a maximalizace entropie. 

%\[ u = \arg\min_{w \in H_{4,6}} \sum_{i =1}^n \frac{|K_i|}{|K|}\log_2|K_i|\]

\begin{veta}[Ekvivalence maximalizace entropie] \label{vetaekvivalencemaxentropy}
    Nechť $K \subset H_{n,k}$ je množina kódů a $u \in H_{n,k}$ je kód. Potom
    \begin{align}\label{rceproentropii}
        \sum_{r\in S_{n,k}, |K_{u,r}| > 0} \frac{|K_{u,r}|}{|K|}\log_{|S_{n,k}|}|K_{u,r}| = \min_{w \in H_{n,k}} \sum_{r\in S_{n,k}, |K_{u,r}| > 0} \frac{|K_{w,r}|}{|K|}\log_{|S_{n,k}|}|K_{w,r}|
    \end{align}
    právě tehdy, když 
    \[h_K(u) = \max_{w \in H_{n,k}} h_K(w).\]
    
\end{veta}
\begin{dukaz} Využíváme toho, že pro $x \in \mathbb{R}$ a konečnou množinu reálných čísel $A$ platí vztahy.
\[z \cdot \min\{a \mid a \in A\} = \min\{z\cdot a \mid a \in A\} \textrm{ pro } z \geq 0\]
a 
\[z \cdot \min \{a \mid a \in A\} = \max\{z\cdot a \mid a \in A\}\textrm{ pro } z < 0.\]
Navíc platí $\log_{|S_{n,k}|} |K_{u,r}| = \frac{\log |K_{u,r}|}{\log |S_{n,k}|}$, a tedy
    \[\log_{2} |K_{u,r}| = \log_{|S_{n,k}|} |K_{u,r}|  \frac{\log |S_{n,k}|}{\log 2}\]
Po přenásobení rovnosti \ref{rceproentropii} konstantou $\frac{\log|S_{n,k}|}{\log2}$ dostáváme
\[\sum_{r\in S_{n,k}, |K_{u,r}| > 0} \frac{|K_{u,r}|}{|K|}\log_{2}|K_{u,r}| = \min_{w \in H_{n,k}} \sum_{r\in S_{n,k}, |K_{w,r}| > 0} \frac{|K_{w,r}|}{|K|}\log_{2}|K_{w,r}|\]
Tuto rovnici přenásobíme konstantou $(-1)$ a převrátíme argument logaritmu.
\[\sum_{r\in S_{n,k}, |K_{u,r}| > 0} \frac{|K_{u,r}|}{|K|}\log_{2}\frac{1}{|K_{u,r}|} = \max_{w \in H_{n,k}} \sum_{r\in S_{n,k}, |K_{w,r}| > 0} \frac{|K_{w,r}|}{|K|}\log_{2}\frac{1}{|K_{w,r}|}\]
Následně k oběma stranám přičteme hodnotu $\log_2|K|$.

\[\sum_{r\in S_{n,k}, |K_{u,r}| > 0} \frac{|K_{u,r}|}{|K|}\log_{2}\frac{|K|}{|K_{u,r}|} = \max_{w \in H_{n,k}} \sum_{r\in S_{n,k}, |K_{w,r}| > 0} \frac{|K_{w,r}|}{|K|}\log_{2}\frac{|K|}{|K_{w,r}|}\]
Využili jsme toho, že $\sum_{i=1}^n \frac{|K^u_i|}{|K|} = 1$ a součtového vzorce pro logaritmus. Z definice $h_K$ plyne 
\[h_K(u) = \max_{w \in H_{n,k}}  h_K(w) .\]
Všechny úpravy byly ekvivalentní, a tedy jsme hotovi.
\end{dukaz}

Algoritmus Max entropy vybírá jako další pokus ten kód, který maximalizuje entropii potomků z definice \ref{defentropierozdeleni}. Dokázali jsme, že stejný kód minimalizuje očekávaný počet pokusů $G$ potřebný k uhádnutí libovolného tajného kódu (\ref{rceocekavanypocetpokusu}). To je myšlenka algoritmu Max entropy.

% Dokázali jsme tedy, že kódy, které minimalizují očekávaný počet pokusů ve smyslu rovnice \ref{rceocekavanypocetpokusu} jsou právě ty kódy, které maximalizují entropii rozdělení z definice \ref{defentropierozdeleni}. To je myšlenka algoritmu Max entropy. 





% Funkce $\log_2(x)$ je rostoucí, a proto tato hodnota určitým způsobem koresponduje s 
% možná by šlo napsat větu, že tabulka rozdělení plyne pouze z kandidátů a dalšího kódu. 



% Tato strategie lze přeformulovat jako minimalizace (přes všechny kódy) očekávaného počtu otázek ano/ne, které jsou nutné pro určení tajného kódu po zahrání dalšího tahu.

% Nechť máme množinu $K$, definujeme počet otázek ano/ne potřebných k určení jakéhokoli prvku $s \in K$ jako $\log_2(|K|)$. Všimněme si, že tato definice koresponduje s příkladem \ref{prbinarysearch}, v případě, když každému prvku množiny 
%V našem případě ale kvůli pravidlům hry nelze jednoduše definovat očekávaný počet ano/ne otázek na určení nějakého kódu, protože se nemůžeme cíleně ptát: "je na i-té pozici k-tá barva?". Navíc nedostáváme odpovědi ano/ne.


% Lze to tedy formulovat i tak, že tah maximalizuje střední hodnotu informačního obsahu po ohodnocení tahu. Informační obsah představuje $-log(\frac{|K_i|}{|K|})$, kde |K| je počet kandidátů před zahráním tahu. Tedy je to nějaká hodnota informace vztažená na počet kandidátů. Algoritmus tedy hledá 


\section{Počet potomků}
Další strategie, kterou navrhl Barteld Kooi \cite{kooi} je volba kódu, který maximalizuje počet ohodnocení, které jako další pokus může dostat. Formálně tento počet definujeme jako počet neprázdných potomků.
\begin{definice}[Počet potomků]
    Nechť $H_{n,k}$ je prostor kódů. Pro $K \subset H_{n,k}$ definujeme valuaci počet potomků jako
    \begin{align*}
        c_K \colon H_{n,k} &\to \mathbb{R} \\
        u &\mapsto |\{r \mid K_{u,r} \neq \emptyset\}|
    \end{align*}
    Prostor všech valuací $c_K, \hspace{3px} K \subset H_{n,k}$ označíme $\mathcal{C}$. 
\end{definice}

\begin{definice}[Algoritmus Most parts]
    \emph{Algoritmus Most parts} definujeme jako deterministický algoritmus s prostorem valuací $\mathcal{C}$ a strategií maximum $M$. Tedy je to funkce, která tajnému kódu $v$ přiřadí posloupnost zahraných kódů s ohodnoceními $\textsc{Solve}[n, k, \mathcal{C}, M](v)$.
\end{definice}


% \begin{definice}[Algoritmus Most parts]
%     % Nechť $n\in \N, k\in \N$. Algoritmus s počtem potomků je pro vstupní tajný kód $v\in H_{n,k}$ definovaný jako $\text{\textsc{SolveMastermind}} (n,k,s_v, \mathcal{C}, M)$. Zde $s_v$ značí zobrazení z definice \ref{defohodnocenivzhledem} a $M$ je funkcionál maximum.

%      Nechť $n\in \N, k\in \N$, $\mathcal{C}$ je prostor valuací počtu potomků a $M$ je strategie maximum. Algoritmus Most parts je definovaný jako funkce 
%      \begin{align*}
%          \textsc{Solve}[n, k, \mathcal{C}, M] \colon H_{n,k} &\to \mathbb{N} \\
%          v & \mapsto \textsc{Solve}[n, k, \mathcal{C}, M](v)
%      \end{align*}
%      podle předpisu algoritmu \ref{alg-default}.
% \end{definice}

\begin{tvrz}[Správnost algoritmu Most parts]
    Algoritmus Most parts řeší [n,k]-Mastermind.
\end{tvrz}
\begin{dukaz}
    Dokážeme, že v každém kroku algoritmu se zmenší velikost množiny kandidátů aktuálního stavu. Platí, že $c_K(u) = 1$ právě tehdy, když $K = K_{u,r}$ pro nějaké $r \in S$. Nechť $K$ množina kandidátů obsahující alespoň dva kódy. Nechť $u \in K$. Potom 
    \[c_K(u) = |\{r \mid K_{u,r} \neq \emptyset\}| \geq 2\]
    protože $|K_{u,(n,0)}| = 1$, a tedy existuje nějaké další $r \in S$, pro které $K_{u,r}$ je neprázdná. V každém kroku se tedy velikost množiny kandidátů aktuálního stavu snižuje. Pokud $|K| = 1$, $c_K(u) = 1$ pro všechny kódy $u \in H_{n,k}$. Algoritmus tedy zvolí jako další pokus jedinný zbývající kód z množiny kandidátů. Protože tajný kód v každém stavu náleží do množiny kandidátů, tak algoritmus tento kód nalezne a hra skončí.
\end{dukaz}


Tato strategie je motivována následujícími příklady, které Kooi \cite{kooi} popisuje ve svém článku. 

\begin{prikl}\label{prdvecasti}
    Uvažujme klasický balíček $52$ karet. Naším úkolem je uhodnout tajnou kartu s tím, že každá karta může být tou tajnou se stejnou pravděpodobností. Předtím se ale můžeme zeptat na jednu ano/ne otázku, kterou snížíme počet zbývajících možností. Jak máme tuto otázku zvolit? 
    Pokud bychom se zeptali, zda je karta srdcová, pravděpodobnost, že kartu po této otázce uhodneme je 
    \[\frac{13}{52} \cdot \frac{1}{13} + \frac{39}{52} \cdot \frac{1}{39} = \frac{2}{52}.\]
    Pokud se zeptáme, zda tajnou kartou je piková sedmička, pravděpodobnost, že kartu po této otázce uhodneme je 
     \[\frac{1}{52} \cdot 1 + \frac{51}{52} \cdot \frac{1}{52} = \frac{2}{52}.\] 
    Nezáleželo tedy, na jak velké množiny nám otázka balíček karet rozdělila. 
\end{prikl}\textbf{}

\begin{prikl}
    Nechť nyní můžeme dostat $k$ různých odpovědí na položenou otázku. Definujeme $K_i$ množinu karet vyhovujících odpovědi $i$. Předpokládáme, že žádná karta nemůže vyhovovat dvěma odpovědím na položenou otázku. Pravděpodobnost, že tajnou kartu uhádneme po této otázce je 
\[\sum_{i = 1}^k \frac{|K_i|}{52} \cdot \frac{1}{|K_i|} = \frac{k}{52}.\]
\end{prikl}


Z těchto dvou příkladů tedy plyne, že nezáleží na velikostech množin, na které otázka rozdělí karty, ale záleží na jejich počtu. Obdobný postup lze aplikovat na hru [n,k]-Mastermind. Jako v případě entropie uvažujeme místo otázek a odpovědí pokusy a jejich ohodnocení. Množiny, na které se množina kandidátů na tajný kód rozdělí jsou neprázdní potomci vzhledem k zahranému kódu. 

%V případě hry [n,k]-Mastermind porovnáváme pravděpodobnosti uhodnutí tajného kódu po daném pokusu. Hledáme tedy kód, který maximalizuje počet možných ohodnocení. Tabulka~\ref{tab02} ukazuje počty možných ohodnocení prvních tahů. Tato strategie by vybrala kód 1123, nebo 1234. 

%Kooi ukázal, že za použití této strategie lze dosáhnout očekávaného počtu pokusů $4.373$. Podrobněji jsou výsledky popsány v následující kapitole.

%V závěru Kooi diskutuje nad možnými zlepšeními/atd - šlo by zmínit








%\subsection{Očekávaný počet}
\chapter{Výsledky}
% V následující kapitole porovnáme tři popisované algoritmy. 

\section{Vlastnosti algoritmů}
Při srovnávání algoritmů budeme používat několik kritérií. První hodnotou je maximální počet pokusů, které algoritmus potřebuje k nalezení libovolného tajného kódu $v \in H_{n,k}$. K nalezení maximálního počtu pokusů stačí vyzkoušet chod algoritmu na všech tajných kódech $v \in H_{n,k}$. V sekci se stromem algoritmu jsme ukázali, že tato hodnota je ekvivalentní největší vzdálenosti listu od kořene ve stromě algoritmu ve smyslu počtu hran.

\begin{definice}[Maximální počet pokusů]\label{defmaxpocetpokusu}
    Nechť $T$ je strom algoritmu pro nějaký algoritmus řešící hru [n,k]-Mastermind. Maximální počet pokusů tohoto algoritmu definujeme jako délku nejdelší cesty od kořene $T$ k nějakému listu $T$. 
\end{definice}

Druhý ukazatel, který u algoritmů budeme srovnávat je průměrný počet pokusů na hru. Ten lze opět odvodit ze stromu algoritmu.

\begin{definice}[Průměrný počet pokusů]\label{defprumpocetpokusu}
    Nechť $T = (V,E)$ je strom algoritmu pro nějaký algoritmus řešící hru [n,k]-Mastermind. Pro list $A$ označíme $d(A)$ jeho vzdálenost od kořene ve smyslu počtu hran. Průměrný počet pokusů algoritmu definujeme jako 
    \[\frac{1}{k^n}\sum_{A\in V, A \text{ list}} d(A).\]
\end{definice}

Posledním kritériem, kterého si budeme u algoritmů všímat je výběr prvního tahu. Zajímá nás, jestli první tah, který algoritmus zvolí, je optimální. Tedy zda neexistuje jiný první pokus, který by s následným použitím algoritmu vracel lepší výsledky. Tuto vlastnost pojmenujeme \emph{konzistence}.

V případě [4,6]-Mastermindu budeme konzistenci zkoumat pro pět počátečních kódů $1111$, $1112$, $1122$, $1123$, $1234$. Obhajujeme to tím, že neexistuje kód, který by vytvořil nějakou jinou valuaci v prvním tahu než tyto kódy. To plyne z toho, že každý kód vznikne nějakou substitucí barev a permutací pozic těchto pěti kódů. Tyto operace nezmění velikosti potomků ${H_{4,6}}_{u,r}$, protože jednotlivé pozice a barvy jsou rovnoměrně rozloženy v celém prostoru kódů. A protože tři použité valuace vycházejí právě z velikostí potomků, jejich hodnoty se nezmění. 

% V sekci $4.5$ diskutujeme možný vliv výběru lexikograficky vyššího prvního tahu z důvodu následného výběru lexikograficky minimálních kódů ve všech algoritmech.

V tabulce \ref{tabvaluaceprvnichtahu} jsou zobrazeny valuace pěti zmíněných kódů vzhledem k $H_{4,6}$. Nejlepší valuace prvního pokusu podle zvolené strategie v konkrétním algoritmu je zobrazena tučně. 


\begin{table}[h]
\centering
\begin{tabular}{l l l l l l}
\toprule
kód & 1111 & 1112 & 1122 & 1123 & 1234 \\
\midrule

Maximum potomka & 625 & 317 & \textbf{256} & 276 & 312 \\
Entropie potomků & 1.50 & 2.69 & 2.89 & 3.04 & \textbf{3.06}\\

% [1.4984350761704437, 2.6934339172716406, 2.885102174947102, 3.0436979819559213, 3.0566709153318907]

%vážený průměr - $\sum_{i = 1}^{14} |S_i|  \cdot \pr (S_i)$ & 512.0 & 235.9 & 204.5 & \textbf{185.3} & 188.2 \\

% [511.9799382716049, 235.94907407407408, 204.53549382716048, 185.26851851851853, 188.18981481481484]

Počet potomků & 5 & 11 & 13 & \textbf{14} & \textbf{14} \\

\bottomrule
\multicolumn{6}{l}{\footnotesize \textit{Pozn:}
Entropie je zaokrouhlena na dvě desetinná místa.
%, vážený průměr na jedno.
}

\end{tabular}
\caption{Hodnoty valuace kódů pro první tah}
\label{tabvaluaceprvnichtahu}
\end{table}



\section{Způsob testování}
Výsledky algoritmů byly nalezeny vytvořením stromu algoritmu pro každou metodu. Strom byl sestavován postupně, kdy pro každý stav (vrchol na grafu) byl nalezen nejlepší další pokus podle zvolené valuace a strategie. Ve chvíli, kdy program dojde do listu hranou $(n,0)$, je uložena aktuální vzdálenost listu od kořene (počet pokusů potřebný k nalezení daného tajného kódu) do seznamu s počty pokusů. 
%V průběhu vytváření stromu byly zjišťovány vzdálenosti konečných stavů (listů) od kořene odpovídající počtu pokusů potřebných k uhádnutí daného kódu. 
Následně z těchto hodnot byly vypočten maximální a průměrný počet pokusů. Konzistence algoritmů byla nalezena sestrojením stromu algoritmu s pevně zvoleným prvním tahem.

Tento postup byl implementován v jazyce Python pro účel této práce. Zdrojový kód a jeho dokumentace je dostupný na GitHub \cite{Simsa_Strategies_for_Mastermind_2025}. 



% Výsledky algoritmů byly nalezeny pomocí programu napsaného v jazyce Python pro účel této práce \cite{Simsa_Strategies_for_Mastermind_2025}. Maximální a průměrné počty pokusů byly nalezeny za pomocí sestavení stromu řešení [4,6]-Mastermind jako na obrázku \ref{fig22minmax}. Strom byl sestavován postupně, kdy pro každý stav (vrchol na grafu) byl nalezen nejlepší další pokus. Ve chvíli, kdy program dojde do listu hranou $(n,0)$, uchová si aktuální počet pokusů potřebný k nalezení daného tajného kódu. Program končí ve chvíli, kdy je každý kód v listu. Nakonec se spočítá očekávaný a maximální počet pokusů ze zachovaných hodnot pro jednotlivé kódy.

% Existují dvě situace, pro které se stav dostane do listu. První je v případě, kdy stav je složen z více kódů a tedy nevíme, který z nich je tajný kódem. V případě, že algoritmus vybere jako další pokus kandidáta, může se stát, že tento kandidát byl dokonce tajným kódem. V tu chvíli máme štěstí a nalezli jsme počet pokusů, který potřebuje min-max algoritmus k nalezení tohoto tajného kódu. Druhá situace nastává, kdy stavem je jednoprvková množina s pouze jedním kódem. Zároveň ho min-max algoritmus zatím nezahrál. V tuto chvíli algoritmus ví, který kód je tajným, ale musí ho ještě zahrát, protože hra končí ve chvíli, kdy je zahraný tajný kód. 

\section{Výsledky algoritmů}

\subsubsection{Min-max}

Tabulka \ref{tabminmaxvysl} ukazuje výsledky Knuthova algoritmu v závislosti na volbě prvního pokusu. Tabulka \ref{tabvaluaceprvnichtahu} ukazuje valuace pěti použitých kódů. Min-max algoritmus tedy jako první tah volí kód $1122$. Maximální počet pokusů Min-max algoritmu je $5$ a tedy je v tomto směru optimální. Průměrně pak potřebuje $4.476$ pokusů. Při volbě jiného prvního tahu jsou ve všech případech oba výsledky horší. Min-max algoritmus je tedy konzistentní.

% Maximální počet pokusů, který min-max algoritmus (začínající kódem 1122) potřebuje na prolomení tajného kódu je $5$. Průměrně pak potřebuje $4.476$ pokusů. K porovnání jsou uvedeny maximální a průměrné počty pokusů pro jiné počáteční kódy. Mohli bychom totiž očekávat, že by průměrný nebo maximální počet pokusů byl menší pro jiný počáteční kód, jak bylo znázorněno v příkladu \ref{malyprikladminmax}. Nejlepších výsledků ale algoritmus dosahuje v při volbě prvního tahu $1122$, který odpovídá přirozenému výběru prvního tahu podle definice Knuthova algoritmu. 

\begin{table}[h]
\centering
\begin{tabular}{l l l l l l}
\toprule
první pokus & 1111 & 1112 & \textbf{1122} & 1123 & 1234 \\
\midrule

maximální počet pokusů 
& 6 & 6 & \textbf{5} & 6 & 6 \\

průměrný počet pokusů 
& 5.065 & 4.556 & \textbf{4.476} & 4.4776 & 4.4776\\
\bottomrule
\multicolumn{6}{l}{\footnotesize \textit{Pozn:}
Průměrný počet pokusů je zaokrouhlen na $3$ desetinná místa}
\end{tabular}
\caption{Výsledky Min-max algoritmu}\label{tabminmaxvysl}
\end{table}

\subsubsection{Max entropy}

Algoritmus Max entropy vybírá jako první pokus kód $1234$, protože tento kód dosahuje největší valuace entropie, jak znázorňuje tabulka \ref{tabvaluaceprvnichtahu}. Tabulka \ref{tabentropievysl} zobrazuje výsledky algoritmu Max entropy. Dosahuje průměrného počtu pokusů $4.415$, což je lepší výsledek než v případě Min-max algoritmu. Tento algoritmus ale nedokáže zaručit vítězství na pět pokusů. 

Navíc, pokud by tento algoritmus volil počáteční kód $1123$, dosáhl by lepšího průměrného počtu pokusů, konkrétně $4.383$. Algoritmus Max entropy tedy není konzistentní. Vysvětlením by mohl být fakt, že entropie neodhaduje ideálně počet pokusů. Mohlo by to být tím, že rozdělení množiny $K$ do potomků $K_{u,r}$ není rovnoměrné. Entropie tedy nedokáže správně odhadnout rozložení množin kandidátů o dva tahy dál. 
%Kdybychom uvažovali valuace, které by srovnávali velikosti potomků do hloubky $2$. 
\begin{table}[h]
\centering
\begin{tabular}{l l l l l l}
\toprule
první pokus & 1111 & 1112 & 1122 & 1123 & \textbf{1234} \\
\midrule

maximální počet pokusů 
& 6 & 6 & 6 & 6 & 6 \\

průměrný počet pokusů 
& 4.911 & 4.496 & 4.428 & \textbf{4.383} & 4.415 \\
\bottomrule
\multicolumn{6}{l}{\footnotesize \textit{Pozn:}
Průměrný počet pokusů je zaokrouhlen na $3$ desetinná místa}
\end{tabular}
\caption{Výsledky algoritmu Max entropy}\label{tabentropievysl}
\end{table}



\subsubsection{Most parts}
Algoritmus Most parts volí jako první tah kód $1123$ (plyne z tabulky \ref{tabvaluaceprvnichtahu}) a dosahuje průměrného počtu pokusů $4.373$. Tento výsledek, zobrazený v tabulce \ref{tabcastivysl}, je ze všech porovnávaných algoritmů nejlepší. Maximální počet pokusů je ale opět horší než u algoritmu Min-max, a to 6. Most parts algoritmus je konzistentní, protože průměrné počty pokusů při odlišném prvním tahu jsou vyšší.


\begin{table}[h]
\centering
\begin{tabular}{l l l l l l}
\toprule
první pokus & 1111 & 1112 & 1122 & \textbf{1123} & 1234 \\
\midrule

maximální počet pokusů 
& 7 & 6 & 6 & 6 & 6 \\

průměrný počet pokusů 
& 4.911 & 4.503 & 4.420 & \textbf{4.373} & 4.444\\
\bottomrule
\multicolumn{6}{l}{\footnotesize \textit{Pozn:}
Průměrný počet pokusů je zaokrouhlen na $3$ desetinná místa}
\end{tabular}
\caption{Výsledky algoritmu Most parts}\label{tabcastivysl}
\end{table}


% \subsection{Očekávaný počet}
% Případná čtvrtá strategie...
% \begin{table}[h]
% \centering
% \begin{tabular}{l l l l l l}
% \toprule
% první pokus & 1111 & 1112 & 1122 & 1123 & 1234 \\
% \midrule

% maximální počet pokusů 
% & 6 & 6 & 5 & 6 & 6 \\

% průměrný počet pokusů 
% & 4.911 & 4.529 & 4.453 & 4.394 & 4.434 \\
% \bottomrule
% \multicolumn{6}{l}{\footnotesize \textit{Pozn:}
% Průměrný počet pokusů je zaokrouhlen na $3$ desetinná místa}
% \end{tabular}
% \caption{Výsledky algoritmu s očekávaným počtem kandidátů}\label{tabminmaxvysl}
% \end{table}

\subsubsection{Srovnání algoritmů}
Nejlepší maximální počet pokusů v případě Min-max algoritmu dává smysl, protože algoritmus minimalizuje maximální počet kódů v potomcích aktuální množiny kandidátů. Fakt, že algoritmus Max entropy dosahuje lepšího průměrného počtu pokusů než Min-max odpovídá tomu, že Max entropy vybírá kód, který minimalizuje očekávanou hodnotu počtu pokusů. 

% Výsledek algoritmu Most parts je do jisté míry překvapivý. Most parts algoritmus totiž maximalizuje pravděpodobnost, že tajný kód v dalším kole uhodneme a nijak se nesnaží odhadnout očekávaný počet pokusů. Domnívám se, že 

% Zároveň fakt, že zbylé algoritmy dosahují lepšího průměrného počtu pokusů také odpovídá tomu, že optimalizují nějakou očekávanou hodnotu následujících stavů. V případě algoritmu Max entropy to je očekávaná hodnota počtu otázek. 

\section{Varianty algoritmů}
\subsubsection{Volba pokusů pouze z kandidátů}
Všechny algoritmy uvedené do této chvíle vybíraly následující pokus ze všech kódů $u\in H_{n,k}$. Předpokládejme, že se nacházíme ve stavu $P$ s množinou kandidátů na tajný kód $K$. Pokud zahrajeme kód, který nenáleží do množiny $K$, nemáme šanci tajný kód v tomto kole prolomit. Logickou variantou algoritmů je způsob, který povoluje výběr dalšího pokusu pouze z množiny kandidátů stavu. V základním algoritmu \ref{alg-default} by množina $U$ na řádce $8$ měla tvar
\[U = \{u \in K \mid f_K(u) = F(f_K)\}.\]
Níže uvádíme výsledky algoritmů volící pouze kandidáty aplikované na [4,6]-Mastermind. Výsledky byly dosaženy stejným postupem jako u algoritmů výše, tedy postupným sestrojením stromu algoritmů. 

\begin{table}[h]
\centering
\begin{tabular}{l l l l}
\toprule
první pokus & Min-max & Max entropy & Počet potomků  \\
\midrule

maximální počet pokusů 
& 6 & 6 & 7  \\

průměrný počet pokusů 
& 4.497 & 4.465 & 4.399 \\
\bottomrule
\multicolumn{4}{l}{\footnotesize \textit{Pozn:}
Průměrný počet pokusů je zaokrouhlen na $3$ desetinná místa}
\end{tabular}
\caption{Výsledky algoritmů volící pouze kandidáty}\label{tabminmaxvysl}
\end{table}

Průměrný počet pokusů se zhoršil ve všech algoritmech, v algoritmu min-max a Max parts ze zhoršil i maximální počet pokusů. Tato volba algoritmů ale zmenší časovou složitost algoritmů. To plyne z toho, že algoritmus v každém stavu algoritmus neporovnává všechny kódy ale pouze kandidáty. 

% \begin{tvrz}
%     Uvažujme algoritmus \ref{alg-default} s výběrem množiny $U$ na řádku $7$ podle následujícího předpisu.
%     \[U = \{u \in K \mid f_K(u) = F(f_K)\}\]
%     Potom tento algoritmus má časovou složitost 
%      \[O \left( \sum_{j = 1}^z  m_j^2 \cdot n^2 \cdot \frac{n^2 + 3n}{2}\right)\]
% \end{tvrz}
% \begin{dukaz}
%     Tvrzení plyne z tvrzení ...
% \end{dukaz}

\subsubsection{Vícekrokové algoritmy}
Všechny popisované algoritmy používaly valuace, které zkoumají rozložení kandidátů u potomků daného stavu. Zkoumali tedy stavy o maximálně jednu hranu v grafu [n,k]-Mastermindu dál. Intuice napovídá, že pokud bychom zkoumali stavy o dvě hrany dál, algoritmy by měli lepší výsledky. Časová složitost těchto algoritmů by ale prudce rostla. Vzhledem k tomu, že optimální strategie pro malé hodnoty $n$ a $k$ již byly nalezeny, tak jsme tyto algoritmy dále nezkoumali. 

% šlo by ale naimplementovat a uvést k porovnání.



\subsubsection{Jiná volba prvního tahu}



\section{Diskuze}
Poznámky:

\begin{itemize}
    \item Algoritmy, které by zkoumaly tabulky rozdělení dva kroky dopředu by jistě byly lepší, výpočetně by to ale bylo mnohem náročnější.
    \item šlo by najít nějakou variantu zmíněných algoritmů, která by řešila mastermind na 5 tahů s výrazně lepším průměrem? - Knuth něco na tenhle způsob ukazoval, ale možná jen v případě jeho algoritmu.
    \item Volba prvního tahu, který by obsahoval vyšší barvy by mohla být lepší, protože algoritmus volí lexikograficky menší kódy. Podobně by bylo možné náhodně promíchat seřazení všech kódů. 
    \item Čím je strategie optimálnější, tím hůře jde popsat.
    \item Strategii optimalizující průměrný počet pokusů nalezl v roce 1993 K. Koyama \cite{koyama}. Ukázal, že lze dosáhnout strategie s průměrným počtem pokusů $4.340$.
    \item aplikace problému mastermind jsou diskutovány v článku doerr playing mastermind with many colours na straně 4
    \item odhad s entropií nefunguje tak dobře. Mohlo by to být silně nerovnoměrně rozdělenými kódy. Entropie neodpovídá ohodnocení, není to binární.

\end{itemize}





\chapter{Poznámky}

\section{Obecné/tipy}
Nějaké tvrzení o n-koulích, např jedno, co je níže

\section{Berghman, Lotte (2008) \cite{BERGHMAN20091880}}
Článek nejprve popisuje existující algoritmy. Jsou řazeny do tří kategorií, Full enumeration, prozkoumání všech kódů a vybrání následujícího pokusu podle zvolené strategie. Tyto algoritmy excelují v průměrném počtu pokusů na hru. Nevýhodou ale je extrémní časová náročnost při procházení všech možných pokračování. 

Druhým typem algoritmů zmíněných je Rules of thumb s cílem dosáhnutí dobrých výsledků za krátký čas. 

Třetím zmíněným typem algoritmu je Meta-heuristics. Tyto algoritmy používají při výběru následujícího pokusu složitější heuristiky, která by měla dosahovat požadovaných výsledků s přihlédnutím na časovou náročnost. 

Autoři popisují svůj program. Jde o genetický algoritmus. Populace je sada možných následujících pokusů. Populace vznikne z předešlé pomocí crossoveru (1-point or 2-point) s pravděpodobností 0.5, mutací, permutací, inverzí. Případně nahrazením vzniklého kódu za náhodný. Eligible codes are added to the list Ei. Výběr z Ei - popsány dva způsoby. První - porovnání s ostatními kódy v Ei - Kolik černých a bílých ohodnocení bychom dostaly při výběru kódu c za předpokladu, že cstar z Ei je tajný kód. Suma přes všěechny cstar z Ei. Výběr buď náhodně, největší číslo nebo nejmenší číslo. Druhý - Vybrat náhodný subset S, Pro každý kód c a každý možný kód cstar je vypočítán počet možných následujících eligible kódů (musíme spočítat body Xc a Yc abychom mohli vybírat eligible kódy). Kód c minimalizující tento součet bude vybrán. Je ukázáno že druhý způsob výběru je efektivnější. 

Nakonec je zmíněný algoritmus porovnán s ostatními algoritmy. Výsledky jsou kvalitní pro rostoucí P a N s výrazně lepším časem. 

Algoritmus vždy začíná tahem 1123, pro který autoři dostávali nejlepší průměrné výsledky. Říkají, že zvolením 1123 jako počáteční tah se ušetří v průměru 0.02 pokusů na prolomení tajného kódu.



\section{Donald E. Knuth (1977) \cite{donald_e__knuth_1977}}

Článek The Computer as Master Mind od Donalda E. Knutha je poměrně krátký, popisuje jeho metodu řešení Logiku na maximálně šest tahů nehledě na tajný kód. Jeho algoritmus se snaží minimalizovat maximální počet zbývajících kandidátů na tajný kód. 

Začíná vždy tahem 1122, pro toto rozložení to funguje. Porovnává to s 1123, který ale nezaručuje vyřešení na pět tahů. V určitých případech může selhat. 

\section{Barteld Kooi - YET ANOTHER MASTERMIND STRATEGY }

Článek popisuje již známé algoritmy/postupy řešení hry mastermind,
popisuje metodu za použití entropie - maximalizuje entropii $-\sum_{i=1}^n p_i log(p_i)$, kde při rozdělení V1 až Vn množiny A je $p_i$ pravděpodobnost, že element z A náleží do Vi. 

Entropie v tomto případě ... maximum entropie odpovídá minimálnímu součtu očekávaného počtu pokusů v rozdělení V1, ... Vn s váhou na velikost těchto rozdělení. 

Je to tam popsáno pěkně - lze se k tomu vrátit.

Dále článek navrhuje nový přístup řešení masterminda. Mluví o tzv strategii rozdělení na nejvíc částí. Jako úvod je nastíněn příklad tipování karty z balíku, kde na pravděpodobnosti uhádnutí nezáleží na velikostech jednotlivých rozdělení ale na jejich počtu. Čím vyšší počet rozdělení balíku, tím větší pravděpodobnost uhádnutí karty v dalším tahu. Toto je zobecněno na pravděpodobnost uhádnutí prvku z A po r+1 otázkách, které postupně rozdělují množinu A na podmnožiny. Díky indukčnímu předpokladu víme, že pravděpodobnost uhádnutí po r+1 otázkách je počet podmnožin na které lze množinu A rozdělit po r+1 otázkách vyděleno velikostí množiny A. 


\section{The Mastermind game and the rigidity of the Hamming space, 2000 \cite{866673}}
Používám definice ze zníměného článku.


\begin{definice}[Hammingův prostor]\label{def01:1}
  Nechť $n \in N$, $q \in N$. Potom množinu všech řetězců (slov) délky $n$ obsahující čísla z množiny $\{1, 2, 3, \dots, q\}$ nazýváme n-dimenzionální q-ární Hammingův prostor a značíme ho $H_{n,q}$.
\end{definice}

\begin{definice}[Hammingova vzdálenost]\label{def01:1}
  Nechť $H_{n,q}$ je Hammingův prostor. Nechť $u,v \in H_{n,q}$. Počet pozic, na kterých se slova $u$ a $v$ liší nazýváme Hammingovou vzdáleností $u$ a $v$, značíme $d(u,v)$.
\end{definice}

\begin{definice}[Báze Hammingova prostoru]\label{def01:1}
  Nechť $H_{n,q}$ je Hammingův prostor. Nechť $B \subset H_{n,q}$. Řekneme, že $B$ je báze Hammingova prostoru, pokud každý prvek z $H_{n,q}$ lze jednoznačně vyjádřit pomocí Hammingových vzdáleností od prvků z $B$.
  \end{definice}

\begin{definice}[Rigidita Hammingova prostoru]\label{def01:1}
  Nechť $H_{n,q}$ je Hammingův prostor. Rigidita Hammingova prostoru, značíme $r_{n,q}$ je minimální možný počet prvků v bázi.
\end{definice}

\begin{veta}\label{veta01:1}
  Nechť $H_{n,q}$ je Hammingův prostor. Potom \[r_{n,q} \geq \frac{n}{log_q(n+1)}\].
\end{veta}

\begin{dukaz}
  Protože počet možných vzdáleností je menší nebo rovno $n+1$ platí, 
  \[(n+1)^{r_{n,q}} \geq q^n\]
  zlogaritmovat
  \[r_{n,q}*log_q(n+1) \geq n\]
  vydelit
  \[r_{n,q} \geq \frac{n}{log_q(n+1)}\].
\end{dukaz}


Z jinych zdroju plati:
$r_{n,q} \geq 2* \frac{n}{log_qn} * (1+o(1))$ 

\[r_{n,2} = 2* \frac{n}{log_2 n} * (1+o(1))\] 



\section{On the metrical rigidity of binary codes () \cite{AVGUSTINOVICH2001444}}
Článek studuje pojem metrical rigidity binárních kódů. Řekneme, že kód je metrically rigid, pokud všechny isomorfismy v metrickém prostoru $C \to E_n$, dále jako isometriky, jdou rozšířit na isometriku celého prostoru $E_n$. 

Kódy jsou ekvivalentní, pokud existuje isometrie, která jeden kód zobrazí na druhý. 


Komentář pro mne:

Když tedy mám neznámý kód se známými vzdálenostmi od nějaké množiny vektorů, potom v každém bude existovat právě jeden????? vektor

Nechť mám kód, který opravuje jednu chybu, má distanci 3. Pokud víme, že kód je metrically rigid, potom pro přijatý vektor s maximálně jednou chybou víme, že jeho distance od nějakého vektoru z kódu je jedna. Navíc věechny jeho distance k ostatním vektorům kódu jsou jasně dané.

Nechť máme bázi metrického prostoru s hammingovou metrikou. potom pro kazdy vektor z prostoru existuje prave jedna kombinace vektorů v bázi, která dá náš vektor. 

Nechť máme dva různé vektory se stejnou hodno

pokud mám r-perfektní redukovaný (obsahuje O) lineární kód, tak pro každý vektor z prostoru je min vzdáleností od kódových slov menší nebo rovna r. Tedy pro 1-perfektní kód je vzdálenost neznámého vektoru y od nějakého kódového slova u rovna buď 0 nebo 1. Pokud je rovna 0, máme vyhráno. Pokud je rovna 1, máme n možností na kód, kde n je dimenze prostoru. Navíc z přítomnosti nuly v kódu známe počet jedniček a nul. Tedy máme buď w(u) nebo n-w(u) možností, kterou pozici změnit podle toho, jestli je vzdálenost y od nuly menší než w(u) nebo větší než w(u). 

\subsection{testovani}
Báze Hammingova [15,11] kódu není báze Hammingova 15, 2 prostoru.
Báze Hammingova [7,4] kódu není báze Hammingova 7, 2 prostoru.


\section{possible idea}
Italo Dejter, Dimacs and University of Puerto Rico

Title: Ternary perfect domination of binary codes

Given a graph G, we say that a vertex subset S of G is a perfect dominating set (or PDS) of G if every vertex v of G not in S is neighbor of exactly one vertex of S.

P. Weichsel showed that each connected component of the subgraph induced by a PDS in the n-cube is a subcube and conjectured that S is uniform, i.e. all such components are isomorphic.

P. Ostergard and D. Weakley showed that the conjecture is false, by exposing a PDS S in the 13-cube whose induced subgraph has as connected components 26 4-cubes and 288 isolated vertices (which correspond to the nonzero words of the ternary Hamming code) and whose automorphism group is the general linear group GL(3,3).

We establish that the perfect dominating set of the 13-cube found by Ostergard and Weakley (as a counterexample to Weichsel's uniformity conjecture) exists very naturally in the ternary Hamming code of length 13, and has a very special geometric-combinatorial structure. This approach provides very short proofs of its characteristic properties

This is joint work with Kevin T. Phelps.

\section{Ternary Hamming and Binary Perfect Covering Codes}
S is a perfect dominating set (or PDS) if for all v in G without S is a neighbour of exactly one vertex of S

Each component of S is an i-cube, where 0 <= i <= n.

Xi is the number of i-cubes in S.
A sphere of radius 1 about any vertex in such i-cube will cover/dominate exactly n-i vertices not in this i-cube. 

Proposition: For a set S to be a PDS the distance between any two components is at least 3 and the number of i-cube components satisfies the bound lower with equality.
\[
\sum_{i = 0}^n 2^i(n-i+1)X_i \leq 2^n
\]

The i/cube components of a PDS can be represented bz codewords over a ternary alphabet. The vertices of an i-cube component all agree on n-i coordinates.

I could use it probably but the construction is problematical for me right now. I would get n-i options in case of an element of the base that is part of an i-cube. I would have to choose between those n-i options which could maybe work.

\section{Testing MDS codes as a base}
\subsection{5-ary [4,2] RS code with alpha = 2}
G = [[1,2,4,3],[1,4,1,4]]
Tenhle kod (25 kodovych slov) z celkem 625 celkovych slov uspel v testu, jestli vsechny tajne kody mohu vyjadrit po techto pokusech. 
Použil jsem hodnocení na black and white pegs

\section{Testing Hamming codes as a base}
\subsection{[7,4] Hamming code}
Hammingův 7,4 kód uspel v testu, jestli vsechny tajne kody mohu vyjadrit po techto pokusech. 
Použil jsem hodnocení na black and white pegs.
Hamming 7,4 má 16 slov, prostor má celkem 128 slov


\section{On the extendability of code isometries}
On the extendability of code isometries is a useful article about proving metrical rigidity of various codes. 

\section{binary 1-perfect code as a base}
If i changed the codeword in a position where the codeword u is same as codeword of distance two from the secret code, the codewords would suddenly be of distance four. Thus, i can only change the codeword u in the positions that are different with the codewords of distance two of the secret code. I have z = (n-1)/2 codewords of distance two from the secret code. If n >= 7, then z >= 3. These codewords have distance three from codeword u. 

If all three codewords are different (wrt u) on the same three positions, then its a contradiction, because they would be the same codewords. If there are two positions where the codewords would all be different from u, then the distance between the codewords would be 2 at most. Therefore there exists maximum of 1 position, where the codewords are all different from u. If there was none, then if we changed anything in u, there would be a codeword with the distance 4 from the resulting word. Thus there exists always one index, where all codewords with distance two from the secret code are different than the codeword u. If we change this index, we get the secret code. 

This concludes the proof, that 1-perfect code suits as a base to static binary mastermind. 

The algorithm stands on finding the index, where all codewords of distance two from the secret code are different than the codeword u. 

So i sum the codewords with codeword u modulus two - i get positions, where the codewords are different. To find the intersection, I pointwise multiply all results.

\section{binary 3-perfect code as a base}
Let C be a 3-perfect binary code. Therefore, d(C) = 7. 7 <= n-k+1. n>=7+k+1. It holds that 3-neighbourhoods of all codewords are disjoint. Also every word is in exactly one 3-neighborhood of some codeword. It means that 3-neighborhood of all words contains exactly one codeword. 

Let v be a secret code. Then there exists a codeword u such that d(u,v) <= 3.

1) If d(u,v) = 0: u = v

2) If d(u,v) = 1: 
There exist n options for index change in u in order to get v. There exist k codewords with distance 6 from v. I have n-1 options for indexes to change in order to get them. 
If there are two options for index change - that means if all k codewords are different than u in these two indexes. Then the distance between the codewords would be 8 at most. I would like to prove that then there would exist a pair of codewords whose distance would be less than 7 which would contradict the definition of 3-perfect binary code.

Maybe the amount of codewords with distance y from secret code is equal to smth*(the amount of codewords with distance y-1 of length n-1 or smth) - find and prove. 

For every word of distance 3 from our secret code v there exists exactly one codeword of distance 3 from our secret code. The number of (feasible) words of distance 3 from our secret code is (n-1) choose 3. Then for every word of these (n-1) choose 3 there exists exactly one codeword of distance three. Some of them are the same codewords though. Let w is a word, z is a codeword of distance three from w. Then there exists 6 choose 3?? words of distance 3 from v such that their distance from z is 3. (I need to prove that)

This might imply that we have ((n-1) choose 3)/(6 choose 3) codewords of distance 6 from the secret code v. 

Now if there existed two indeces i,j such that. u + ei is a secret code and u + ej is a secret code. Therefore all ((n-1) choose 3)/(6 choose 3) are the same on those two indeces. I want to imply that there exists a pair of codewords out of these that are equal on three indeces, therefore their distance is at most 6, which is a contradiction. There is n-2 options for the third index. We want all codewords that are different on four positions out of the n-2 options. Therefore ((n-1) choose 3)/(6 choose 3) * 4 has to be less than or equal to n-2. 

(Pocet kodovych slov krat pocet indexu (ten zlomek), na kterych musi byt odlisne (4)) chci mit vetsi nez pocet indexů, na kterych mohou byt odlisne (n-2 atd). Potom vime, ze nejaka dve kodova slova se shoduji na vice pozicich, nez mohou, a tim dojdu ke sporu s minimální moznou vzdalenosti mezi kodovymi slovy. 

V rovnicích chybí LS * 4.

\[
\frac{\binom{n-1}{3}}{\binom{6}{3}} = \frac{\frac{(n-1)*(n-2)*(n-3)}{6}}{\frac{6*5*4}{6}} = \frac{(n-1)*(n-2)*(n-3)}{6*5*4}
\]
So we need 
\[
\frac{(n-1)*(n-2)*(n-3)}{6*5*4} > n-2
\]

\[(n-1)*(n-3) > 6*5*4\]

\[
n^2-4n+3 > 120
\]

\[
n^2-4n-117 > 0
\]
which holds for n < -9, n > 13. 

We want n > 13. 

Construct similar proof for distance 2 and 3 from the closest codeword:

Let d(u,v) = 2:

There exist $\binom{n-2}{2}$ words of distance two from v. There exist $\frac{\binom{n-2}{2}}{\binom{5}{2}}$
We want 
\[
\frac{\binom{n-2}{2}}{\binom{5}{2}} > n-3
\]
\[\frac{(n-2)*(n-3)}{20} > n-3\]
\[n-2 > 20\]
\[n > 22\]

Let d(u,v) = 3:
\[\frac{\binom{n-3}{1}}{\binom{4}{1}} > n-4\]
\[\frac{n-3}{4} > n-4\]
\[n-3 > 4*n-16\]
\[13 > 3n\]
\[n < 13/3 \approx 4.3\]
\[n \leq 4\]

I would like to use somehow the white pegs. Do they tell us just the number of 1/0 in the secret code? Or do they say something else? I have the weight of the secret code w(v) from the distance from zero. 

If the distance from u to v is 3, then the options for the weight of v are either w(u)-3, w(u)-1, w(u)+1, w(u)+3. 

Let u be wrong in the first three coordinates. 

\begin{itemize}
    \item u \approx 000, then v \approx 111. Then w(v) = w(u) + 3. The number of white pegs is zero. 
    \item u \approx 001, Then v \approx 110. Then w(v) = w(u) + 1. The number of white pegs is two.
    \item u \approx 011, then v \approx 100. Then w(v) = w(u) - 1. The number of white pegs is two. 
    \item u \approx 111, then v \approx 000. Then w(v) = w(u) - 3. The number of white pegs is zero. 
\end{itemize}

Let k be the number of white pegs corresponding to u. Then k is either zero or two. 

If k is zero, we have w(u) choose 3 or (n - w(u)) choose three possibilities. Which option we choose, that is according to the distance from zero (can be determined). If we have $\binom{w(u)}{3}$ options. We change 1->0. Then let there exist $\binom{n-3}{1}$ eligible words (Z) of distance one from v. These have the weight of w(u)-4 or w(u)-2. They are of distance 3 from certain codeword. The number of elements of Z with the same codeword in the neighborhood is 4. I have n-4 options 

Chci pocet indexu, na kterych se kodova slova lisi od u krat pocet 

\section{proposition for r-perfect binary codes}
Proposition: for r in the set of natural numbers there exist n0 such that every r-perfect binary code of length n>=n0 is a base of the hamming space $H_{2,n}$.

Not proven yet but could work maybe.


\chapter*{Závěr}
\addcontentsline{toc}{chapter}{Závěr}


%%% Seznam použité literatury
%%% Seznam použité literatury (bibliografie)
%%%
%%% Pro vytváření bibliografie používáme biblatex. Ten zpracovává
%%% citace v textu (např. makro \cite{...}) a vyhledává k nim literaturu
%%% v souboru literatura.bib.
%%%
%%% Podívejte se na nastavení biblatexu v souboru thesis.tex.

%%% Vytvoření seznamu literatury. Pozor, pokud jste necitovali ani jednu
%%% položku, seznam se automaticky vynechá.

% Dovolíme položkám trochu vyčuhovat přes pravý okraj.
\def\bibfont{\hfuzz=2pt}

\printbibliography[heading=bibintoc,title=Literatura]

%%% Kdybyste chtěli bibliografii vytvářet ručně (bez biblatexu), lze to udělat
%%% následovně. V takovém případě se řiďte normou ISO 690 a zvyklostmi v oboru.

% \begin{thebibliography}{99}
%
% \bibitem{lamport94}
%   {\sc Lamport,} Leslie.
%   \emph{\LaTeX: A Document Preparation System}.
%   2. vydání.
%   Massachusetts: Addison Wesley, 1994.
%   ISBN 0-201-52983-1.
%
% \end{thebibliography}


%%% Obrázky v práci
%%% (pokud jich je malé množství, obvykle není třeba seznam uvádět)
\listoffigures

%%% Tabulky v práci (opět nemusí být nutné uvádět)
%%% U matematických prací může být lepší přemístit seznam tabulek na začátek práce.
\listoftables

%%% Použité zkratky v práci (opět nemusí být nutné uvádět)
%%% U matematických prací může být lepší přemístit seznam zkratek na začátek práce.
\chapwithtoc{Seznam použitých zkratek}

%%% Součástí doktorských prací musí být seznam vlastních publikací
\ifx\ThesisType\TypePhD
\chapwithtoc{Seznam publikací}
\fi

%%% Přílohy k práci, existují-li. Každá příloha musí být alespoň jednou
%%% odkazována z vlastního textu práce. Přílohy se číslují.
%%%
%%% Do tištěné verze se spíše hodí přílohy, které lze číst a prohlížet (dodatečné
%%% tabulky a grafy, různé textové doplňky, ukázky výstupů z počítačových programů,
%%% apod.). Do elektronické verze se hodí přílohy, které budou spíše používány
%%% v elektronické podobě než čteny (zdrojové kódy programů, datové soubory,
%%% interaktivní grafy apod.). Elektronické přílohy se nahrávají do SISu.
%%% Povolené formáty souborů specifikuje opatření rektora č. 72/2017.
%%% Výjimky schvaluje fakultní koordinátor pro zavěrečné práce.
\appendix
\chapter{Přílohy}

\section{První příloha}

\end{document}
